% =============================================================================
%  Tài liệu 03 - Kế hoạch quản lý dự án (Project Management Plan)
%  Build: latexmk 03-project-management-plan.tex   (pdfLaTeX, output in build/)
% =============================================================================
\documentclass[12pt,a4paper]{report}
\usepackage{sumtalk-report}
\usepackage{pgfgantt}
% =============================================================================
%  sumtalk-info.tex - information shared by every document (01-06).
%  Edit here once; every cover and footer picks it up.
% =============================================================================
\newcommand{\ProjectName}{SumTalk}
\newcommand{\CourseName}{QUẢN LÝ DỰ ÁN PHẦN MỀM}          % cover (upper case)
\newcommand{\CourseNameNormal}{Quản lý Dự án Phần mềm}     % footer
\newcommand{\AcademicYear}{2026 -- 2027}
\newcommand{\Instructor}{\canDien[họ tên giảng viên]}

% One \item per member: "Full name -- student ID", sorted by given name.
\newcommand{\TeamMembers}{%
  \item Nguyễn Anh Duy -- \canDien[MSSV]
  \item Đỗ Thành Đạt -- \canDien[MSSV]
  \item Trần Hải Đăng -- \canDien[MSSV]
  \item Châu Hùng Hào -- \canDien[MSSV]
  \item Bùi Đoàn Minh Nhật -- \canDien[MSSV]
  \item Võ Hoàng Sơn -- \canDien[MSSV]
}


\newcommand{\DocTitle}{KẾ HOẠCH QUẢN LÝ DỰ ÁN}
\newcommand{\DocTitleNormal}{Kế hoạch quản lý dự án}
\newcommand{\DocTitleEn}{Project Management Plan}
\newcommand{\DocRunningTitle}{Kế hoạch quản lý dự án}
\newcommand{\DocCode}{SUMTALK-PMP}
\newcommand{\DocVersion}{2.1 -- 24/11/2026}
\newcommand{\DocStatus}{Chờ duyệt}
\newcommand{\DocAuthor}{\canDien[Project Manager]}
\newcommand{\DocReviewers}{Sponsor, Product Owner, Project Manager}
\newcommand{\DocReaders}{Sponsor, Project Manager, cả nhóm dự án}
\newcommand{\CoverDate}{TP. HỒ CHÍ MINH, THÁNG 11/2026}
\hypersetup{pdftitle={SumTalk -- Kế hoạch quản lý dự án},pdfauthor={Nhóm SumTalk}}

% Role placeholder in the risk register (a role that still has to be replaced by a name).
\newcommand{\role}[1]{\textit{#1}}

\begin{document}

\makereportcover

% ============================================================ Front matter
\chapter*{Thông tin tài liệu}
\addcontentsline{toc}{chapter}{Thông tin tài liệu}

\docinfotable

\section*{Lịch sử phiên bản}

\begin{table}[H]
  \centering
  \begin{tabularx}{\linewidth}{@{}Q{1.6cm}Q{2.3cm}LP{3.2cm}@{}}
    \toprule
    \thead{Phiên bản} & \thead{Ngày} & \thead{Nội dung thay đổi} & \thead{Người thực hiện} \\
    \midrule
    1.0 & 12/11/2026 & Bản đầu, gồm bảy tài liệu riêng: WBS, ước lượng effort, lịch trình CPM/Gantt, Risk Register,
      kế hoạch nguồn lực, kế hoạch chi phí, RACI và truyền thông & \canDien \\
    2.0 & 13/11/2026 & Gộp thành một tài liệu, đánh lại số mục (CR-08) & \canDien \\
    2.1 & 24/11/2026 & Chuyển sang định dạng báo cáo PDF (LaTeX), sắp xếp lại thành chín chương; vẽ cây WBS, sơ đồ
      mạng, Gantt, ma trận rủi ro; sửa tổng story point của bảng ước lượng (118 SP) & \canDien \\
    \bottomrule
  \end{tabularx}
\end{table}

\needspace{8\baselineskip}
\heading{Quy ước trong tài liệu}
\begin{itemize}
  \item \canDien[\ldots]: số liệu đã có ở Azure Boards, dashboard chi phí hoặc biên bản họp nhưng chưa chép
        vào tài liệu; cần bổ sung trước khi trình duyệt.
  \item \textbf{Giả định}: con số đặt ra để tính toán; bảng ghi rõ công thức để tính lại khi có số thật, hoặc
        tính lại bằng công cụ \texttt{pm\_calc.py}~\cite{pmcalc}.
  \item \textbf{Thời lượng} trong lịch trình tính bằng \emph{ngày làm việc} và là số ước lượng từ story point.
  \item \textbf{Story point (SP)} là điểm đề xuất để đối chiếu với điểm đã vote trên Azure Boards.
  \item \textbf{Mã công việc}: \#61--\#72, \#99, \#100 là ID trên Azure Boards; T1.x là task của epic E1; A--W là
        mã công việc trong lịch trình (Chương~\ref{chap:schedule}).
  \item Tên \role{in nghiêng} ở cột phụ trách là \emph{vai trò}, cần thay bằng tên thành viên.
\end{itemize}

\tableofcontents

\chapter*{Tóm tắt}
\addcontentsline{toc}{chapter}{Tóm tắt}

Tài liệu này là \textbf{kế hoạch quản lý dự án} (Project Management Plan -- PMP) của SumTalk: kế hoạch tổng thể
để thực hiện và kiểm soát dự án sau khi Sponsor đã duyệt khởi động. Phạm vi công việc được phân rã thành
\textbf{9 nhánh cấp 1 và 33 gói công việc}, mỗi gói tương ứng một task trên Azure Boards hoặc một nhóm commit
trên GitHub. Effort được ước lượng bằng story point với Planning Poker (\textbf{118 SP} đề xuất cho 18 hạng mục) và
được đối chiếu bằng hai phương pháp độc lập: Function Point cho khoảng \textbf{1.560 giờ} phần chức năng,
COCOMO~81 cho \textbf{3.500--5.100 giờ} nếu làm theo kiểu công ty với nhóm toàn thời gian.

Lịch trình gồm 23 công việc; đường găng A $\rightarrow$ B $\rightarrow$ C/F $\rightarrow$ G $\rightarrow$ U
$\rightarrow$ V $\rightarrow$ W dài \textbf{33 ngày làm việc}, còn nhánh tóm tắt và dịch có 3--12 ngày dự trữ.
Risk Register theo dõi \textbf{14 rủi ro}, trong đó 5 rủi ro ở vùng Cao; rủi ro lớn nhất hiện nay là chi phí và
giới hạn của API bên ngoài khi deploy, cùng độ chính xác STT với tiếng Việt và họp trộn ngôn ngữ. Chi phí chủ
yếu là nhân công với đơn giá \textbf{100.000~VND/giờ}; hạ tầng và API cho bản test chỉ khoảng
\textbf{2,5--3 triệu VND/tháng} khi chạy 24/7, ngân sách giữ thêm 10\% dự phòng. Cuối cùng, ma trận RACI giao
mỗi sản phẩm cho đúng một người chịu trách nhiệm cuối, kèm lịch họp Scrum, kế hoạch truyền thông và quy trình
leo thang.

\textbf{Từ khoá:} kế hoạch quản lý dự án, WBS, story point, COCOMO, Function Point, CPM, Gantt, Risk Register,
ngân sách, RACI.

\chapter*{Danh mục từ viết tắt}
\addcontentsline{toc}{chapter}{Danh mục từ viết tắt}

% 39 entries: longtable so the list can continue on the next page.
\begingroup
  \setstretch{1.1}
  \begin{longtable}{@{}P{2.2cm}P{13.2cm}@{}}
    \toprule
    \thead{Viết tắt} & \thead{Ý nghĩa} \\
    \midrule
    \endfirsthead
    \toprule
    \thead{Viết tắt} & \thead{Ý nghĩa} \\
    \midrule
    \endhead
    \bottomrule
    \endlastfoot
    AC     & Actual Cost -- chi phí thực tế \\
    AFP    & Adjusted Function Point -- điểm chức năng đã điều chỉnh \\
    AON    & Activity on Node -- sơ đồ mạng đặt công việc trên nút \\
    API    & Application Programming Interface -- giao diện lập trình ứng dụng \\
    AWS    & Amazon Web Services (EC2: máy chủ ảo; EBS: ổ đĩa; S3: lưu trữ đối tượng) \\
    BA     & Business Analyst -- chuyên viên phân tích nghiệp vụ \\
    BAC    & Budget at Completion -- tổng ngân sách khi hoàn thành \\
    COCOMO & Constructive Cost Model -- mô hình ước lượng chi phí phần mềm \\
    CPI    & Cost Performance Index -- chỉ số hiệu quả chi phí \\
    CPM    & Critical Path Method -- phương pháp đường găng \\
    CPU    & Central Processing Unit -- bộ xử lý trung tâm \\
    E2E    & End-to-End -- kiểm thử đầu cuối \\
    EVA    & Earned Value Analysis -- phân tích giá trị thu được \\
    FP     & Function Point -- điểm chức năng \\
    GPU    & Graphics Processing Unit -- bộ xử lý đồ hoạ, dùng để chạy model \\
    HTTPS  & Hypertext Transfer Protocol Secure -- giao thức web có mã hoá \\
    IAM    & Identity and Access Management -- quản lý danh tính và quyền trên AWS \\
    IFPUG  & International Function Point Users Group -- tổ chức chuẩn hoá cách đếm FP \\
    KLOC   & Kilo Lines of Code -- nghìn dòng mã \\
    LLM    & Large Language Model -- mô hình ngôn ngữ lớn \\
    PM     & Project Manager -- quản lý dự án \\
    PMP    & Project Management Plan -- kế hoạch quản lý dự án \\
    PO     & Product Owner -- chủ sản phẩm \\
    POC    & Proof of Concept -- bản thử nghiệm khái niệm \\
    PV     & Planned Value -- giá trị (chi phí) kế hoạch \\
    PWA    & Progressive Web App -- ứng dụng web cài được như app \\
    QC     & Quality Control -- kiểm soát chất lượng \\
    RACI   & Responsible, Accountable, Consulted, Informed -- ma trận phân vai \\
    SM     & Scrum Master \\
    SP     & Story Point -- điểm ước lượng độ lớn của story \\
    SPI    & Schedule Performance Index -- chỉ số hiệu quả tiến độ \\
    STT    & Speech-to-Text -- chuyển giọng nói thành văn bản \\
    UAT    & User Acceptance Testing -- kiểm thử chấp nhận của người dùng \\
    UFP    & Unadjusted Function Point -- điểm chức năng chưa điều chỉnh \\
    UI/UX  & User Interface / User Experience -- giao diện / trải nghiệm người dùng \\
    VAD    & Voice Activity Detection -- phát hiện đoạn có tiếng nói \\
    VAF    & Value Adjustment Factor -- hệ số điều chỉnh của FP \\
    WBS    & Work Breakdown Structure -- cấu trúc phân rã công việc \\
  \end{longtable}
\endgroup

\cleardoublepage\phantomsection
\addcontentsline{toc}{chapter}{\listfigurename}
\listoffigures

\cleardoublepage\phantomsection
\addcontentsline{toc}{chapter}{\listtablename}
\listoftables

% ============================================================ Chapter 1
\chapter{Giới thiệu}
\label{chap:intro}

\section{Bối cảnh}

Hồ sơ khởi động dự án~\cite{pid} đã xác lập SumTalk, đặt sáu mục tiêu SMART và chọn \textbf{Scrum} làm phương
pháp quản lý, kèm một bộ tài liệu quản lý tối thiểu để báo cáo với Sponsor. Đặc tả yêu cầu phần
mềm~\cite{srs} mô tả sản phẩm cần làm: 7 epic, 17 user story và 37 yêu cầu phi chức năng. Tài liệu này trả lời
câu hỏi tiếp theo: dự án sẽ được \emph{thực hiện} và \emph{kiểm soát} như thế nào để giao đúng sản phẩm đó,
trong thời gian và ngân sách đã duyệt.

\section{Mục đích và người đọc của tài liệu}

Kế hoạch bao gồm bảy lĩnh vực: phạm vi công việc (WBS), ước lượng effort, lịch trình và đường găng, rủi ro,
nguồn lực, chi phí, tổ chức nhóm và truyền thông. Các khái niệm và kỹ thuật dùng ở đây (WBS, CPM, ma trận rủi ro,
RACI) theo hướng dẫn PMBOK~\cite{pmbok}. Kế hoạch chất lượng và kiểm thử không nằm trong tài liệu này mà ở
\docref{04}{Chất lượng và kiểm thử}~\cite{qtd}; số liệu thực tế của từng sprint (burndown, EVA, Change Log) được
báo cáo trong \docref{05}{Báo cáo theo dõi và kiểm soát}~\cite{mcr}.

Người đọc chính là Sponsor (duyệt lịch, ngân sách và mức rủi ro chấp nhận được), Project Manager (dùng kế hoạch
làm đường cơ sở để theo dõi) và cả nhóm dự án (biết việc của mình, phụ thuộc và người cần phối hợp). Theo cách áp
dụng Scrum đã chọn, kế hoạch \emph{được cập nhật mỗi sprint}, không lập một lần rồi khoá.

\section{Cấu trúc tài liệu}

Phần còn lại của tài liệu gồm tám chương:
\begin{itemize}
  \item \textbf{Chương~\ref{chap:wbs} -- Phạm vi công việc (WBS):} cây WBS, từ điển WBS đến gói công việc và quy
        tắc phân rã.
  \item \textbf{Chương~\ref{chap:effort} -- Ước lượng effort:} story point và Planning Poker, velocity, đối chiếu
        bằng COCOMO~81 và Function Point.
  \item \textbf{Chương~\ref{chap:schedule} -- Lịch trình:} danh sách công việc, CPM, sơ đồ mạng, Gantt và release
        plan theo sprint.
  \item \textbf{Chương~\ref{chap:risk} -- Quản lý rủi ro:} thang đánh giá, Risk Register, ma trận rủi ro, quy trình
        và dự phòng.
  \item \textbf{Chương~\ref{chap:resource} -- Nguồn lực:} con người, phân công theo sprint, san tải, máy móc và
        dịch vụ.
  \item \textbf{Chương~\ref{chap:cost} -- Chi phí và ngân sách:} đơn giá, chi phí nhân công và ngoài nhân công,
        tổng ngân sách và cách kiểm soát.
  \item \textbf{Chương~\ref{chap:team} -- Tổ chức nhóm và truyền thông:} vai trò, ma trận RACI, lịch họp Scrum,
        kế hoạch truyền thông và leo thang vấn đề.
  \item \textbf{Chương~\ref{chap:conclusion} -- Kết luận và phê duyệt.}
\end{itemize}

% ============================================================ Chapter 2
\chapter{Phạm vi công việc (WBS)}
\label{chap:wbs}

\textbf{Dự án được phân rã thành 9 nhánh cấp 1 và 33 gói công việc.} Nhánh đầu là quản lý dự án, bảy nhánh tiếp
theo ứng với bảy epic của Product Backlog~\cite{srs}, nhánh cuối là kiểm thử và phát hành. Gói công việc ở cấp
thấp nhất tương ứng với task trên Azure Boards (T1.x, \#61--\#72, \#99, \#100) hoặc các commit theo task trên
GitHub.

Hình~\ref{fig:wbs} trình bày WBS đến cấp 2. Hai nhánh tóm tắt (E3) và dịch (E4) có nhiều gói nhất, mỗi nhánh 5
gói, vì ngoài model, lưu trữ, job và giao diện còn có một gói đánh giá chất lượng (\#72, \#66).

\begin{figure}[htbp]
  \centering
  \begin{tikzpicture}[
      root/.style={draw=primaryblue,fill=primaryblue,text=white,font=\small\bfseries,
                   rounded corners=2pt,inner sep=4pt,anchor=west},
      l1/.style={draw=primaryblue,fill=accentblue,font=\footnotesize,align=left,text width=3.4cm,
                 minimum height=0.85cm,inner sep=3pt,rounded corners=1.5pt,anchor=west,
                 execute at begin node={\setstretch{1}\hyphenpenalty=10000\relax}},
      l2/.style={draw=primaryblue!60,fill=white,font=\scriptsize,align=left,text width=1.9cm,
                 minimum height=0.85cm,inner sep=3pt,anchor=west,
                 execute at begin node={\setstretch{1}\hyphenpenalty=10000\relax}},
      ln/.style={primaryblue!70,thin}]
    \node[root] (root) at (0,0) {SumTalk};
    \draw[ln] (0.35,-0.3) -- (0.35,-9.9);
    \foreach \lab/\kids [count=\i] in {
        {\textbf{1} -- Quản lý dự án}/{1.1 Khởi động,{1.2 Lập kế hoạch},{1.3 Theo dõi, kiểm soát},1.4 Kết thúc},
        {\textbf{2 -- E1} Ghi lại cuộc họp}/{{2.1 Ghi âm trên trình duyệt},{2.2 Offline, điện thoại},2.3 Upload},
        {\textbf{3 -- E2} Transcript}/{{3.1 STT ngôn ngữ gốc},{3.2 Transcript live},{3.3 Trang cuộc họp},3.4 File dài},
        {\textbf{4 -- E3} Tóm tắt (\#67)}/{{4.1 Model, prompt},4.2 Map-reduce,{4.3 Job, lưu trữ, API},{4.4 Giao diện},{4.5 Đánh giá}},
        {\textbf{5 -- E4} Dịch (\#60)}/{5.1 Model dịch,{5.2 Lưu bản dịch, API},{5.3 Job có tiến độ},{5.4 Giao diện, xuất file},5.5 Đánh giá},
        {\textbf{6 -- E5} App desktop}/{6.1 Vỏ Electron,6.2 Đóng gói,{6.3 Thích nghi web app}},
        {\textbf{7 -- E6} Engine STT dùng chung}/{{7.1 Client STT remote},{7.2 API STT chuẩn OpenAI},{7.3 Cấu hình bộ cài}},
        {\textbf{8 -- E7} Triển khai web AWS}/{8.1 Lưu trữ S3,{8.2 Container, proxy},{8.3 Hướng dẫn deploy}},
        {\textbf{9} -- Kiểm thử và phát hành}/{9.1 Test tự động,{9.2 Test E2E Sprint 3},{9.3 UAT, phát hành}}} {
      \pgfmathsetmacro{\y}{-1.1*\i}
      \draw[ln] (0.35,\y) -- (0.6,\y);
      \node[l1] (L\i) at (0.6,\y) {\lab};
      \foreach \k [count=\n] in \kids {
        \pgfmathsetmacro{\xk}{4.5+(\n-1)*2.25}
        \ifnum\n=1 \pgfmathsetmacro{\xs}{4.21}\else\pgfmathsetmacro{\xs}{\xk-0.14}\fi
        \draw[ln] (\xs,\y) -- (\xk,\y);
        \node[l2] at (\xk,\y) {\k};
      }
    }
  \end{tikzpicture}
  \caption{Cấu trúc phân rã công việc (WBS) của dự án SumTalk, cấp 1--2}
  \label{fig:wbs}
\end{figure}

\section{Từ điển WBS}

Từ điển WBS (Bảng~\ref{tab:wbs-dict}) mô tả sản phẩm bàn giao của từng gói công việc và nơi theo dõi gói đó:
tài liệu trong bộ quản lý dự án, task trên Azure Boards, hoặc file và nhánh trong kho mã nguồn.

\begingroup\small
\setlength{\tabcolsep}{4pt}
\begin{longtable}{@{}P{1.0cm}P{3.2cm}P{6.3cm}P{4.4cm}@{}}
  \caption{Từ điển WBS đến gói công việc}\label{tab:wbs-dict}\\
  \toprule
  \thead{Mã} & \thead{Gói công việc} & \thead{Sản phẩm bàn giao} & \thead{Liên kết board / repo} \\
  \midrule
  \endfirsthead
  \multicolumn{4}{@{}l}{\footnotesize\itshape Bảng~\thetable{} (tiếp theo)}\\
  \toprule
  \thead{Mã} & \thead{Gói công việc} & \thead{Sản phẩm bàn giao} & \thead{Liên kết board / repo} \\
  \midrule
  \endhead
  \bottomrule
  \endlastfoot
  \multicolumn{4}{@{}l}{\textbf{1 -- Quản lý dự án}}\\
  1.1 & Khởi động & Charter, Business Case, chọn phương pháp & Hồ sơ khởi động~\cite{pid} \\
  1.2 & Lập kế hoạch & Backlog, yêu cầu, ước lượng, lịch, rủi ro, nguồn lực, chi phí, chất lượng, RACI
      & Đặc tả yêu cầu~\cite{srs}, tài liệu này, Chất lượng và kiểm thử~\cite{qtd} \\
  1.3 & Theo dõi và kiểm soát & Sprint report, EVA, biên bản, Change Log & Báo cáo theo dõi và kiểm soát~\cite{mcr} \\
  1.4 & Kết thúc & SOW và nghiệm thu, test summary, release notes, lessons learned
      & Báo cáo kết thúc~\cite{pcr}, Test Summary~\cite{qtd} \\
  \midrule
  \multicolumn{4}{@{}l}{\textbf{2 -- E1 Ghi lại cuộc họp}}\\
  2.1 & Ghi âm trên trình duyệt & Màn hình ghi âm, recorder (micro + âm thanh họp), gửi chunk, tạo meeting
      & T1.1.1--T1.1.5 \\
  2.2 & Offline và điện thoại & IndexedDB, đồng bộ, Wake Lock, PWA & T1.2.1--T1.2.5 \\
  2.3 & Upload & Màn hình upload, kiểm tra file, upload theo phần, chuẩn hoá audio & T1.3.1--T1.3.5 \\
  \midrule
  \multicolumn{4}{@{}l}{\textbf{3 -- E2 Transcript}}\\
  3.1 & STT ngôn ngữ gốc & Whisper GPU/CPU, nhận ngôn ngữ từng câu, lọc chữ ``ảo'' & \path{services/transcriber.py} \\
  3.2 & Transcript live & Live theo LocalAgreement, chunk 1\,s, chữ chắc chắn và chữ tạm
      & Nhánh \path{feature/streaming-stt} \\
  3.3 & Trang cuộc họp & Player đồng bộ, tìm kiếm, xuất TXT/SRT/VTT/JSON/WAV, phiên âm lại
      & \path{pages/MeetingPage.tsx} \\
  3.4 & File dài & Chia đoạn tại chỗ im lặng, lưu dần & Nhánh \path{feature/streaming-stt} \\
  \midrule
  \multicolumn{4}{@{}l}{\textbf{4 -- E3 Tóm tắt (\#67)}}\\
  4.1 & Model và prompt tóm tắt & Schema, prompt, client LLM & \#68 \\
  4.2 & Transcript dài & Map-reduce & \#69 \\
  4.3 & Job tóm tắt, lưu trữ, API & Worker, retry, model dữ liệu, endpoint & \#70 \\
  4.4 & Giao diện tóm tắt & Tab Summary và Regenerate & \#71 \\
  4.5 & Đánh giá tóm tắt & \path{eval_summary.py} & \#72 \\
  \midrule
  \multicolumn{4}{@{}l}{\textbf{5 -- E4 Dịch (\#60)}}\\
  5.1 & Model dịch & Prompt, batch, nhận ngôn ngữ & \#61 \\
  5.2 & Lưu bản dịch, API & Model dữ liệu, endpoint & \#62 \\
  5.3 & Job dịch có tiến độ & Worker & \#63 \\
  5.4 & Giao diện và xuất & Chọn ngôn ngữ, song ngữ, xuất TXT/SRT & \#64, \#65 \\
  5.5 & Đánh giá dịch & \path{eval_translation.py} & \#66 \\
  \midrule
  \multicolumn{4}{@{}l}{\textbf{6 -- E5 App desktop}}\\
  6.1 & Vỏ Electron & Cửa sổ, menu, quyền, vòng đời backend & \path{desktop/main.js}, \path{lib/backend.js} \\
  6.2 & Đóng gói & PyInstaller, bộ cài NSIS, bản GPU
      & \path{desktop/sumtalk-server.spec}, \path{build-desktop.bat} \\
  6.3 & Thích nghi web app & Nhận biết desktop, âm thanh hệ thống & \path{lib/desktop.ts} \\
  \midrule
  \multicolumn{4}{@{}l}{\textbf{7 -- E6 Engine STT dùng chung}}\\
  7.1 & Client STT remote & VAD, chia block, song song, retry & \path{services/stt_remote.py} \\
  7.2 & API STT chuẩn OpenAI & \path{/v1/audio/transcriptions}, API key & \path{api/stt_server.py} \\
  7.3 & Cấu hình bộ cài & \texttt{SUMTALK\_DEFAULT\_*} & \path{desktop/scripts/build-backend.js} \\
  \midrule
  \multicolumn{4}{@{}l}{\textbf{8 -- E7 Triển khai web AWS}}\\
  8.1 & Lưu trữ S3 & Lưu audio, presigned URL, sao lưu database & \path{services/storage.py} \\
  8.2 & Container và proxy & Dockerfile, Compose, Caddy & \path{Dockerfile}, \path{deploy/} \\
  8.3 & Hướng dẫn deploy & Tài liệu từng bước & \path{docs/deploy/AWS-EC2-S3.md} \\
  \midrule
  \multicolumn{4}{@{}l}{\textbf{9 -- Kiểm thử và phát hành}}\\
  9.1 & Test tự động & pytest, Playwright (trình duyệt, desktop)
      & \path{backend/tests}, \path{frontend/e2e}, \path{desktop/scripts/smoke-test.js} \\
  9.2 & Test E2E Sprint 3 & Test case, kịch bản & \#99, \#100 \\
  9.3 & UAT và phát hành & Biên bản UAT, release notes, bộ cài, bản web
      & Test Summary Report~\cite{qtd}, Release Notes~\cite{pcr} \\
\end{longtable}
\endgroup

\section{Quy tắc phân rã}

Ba quy tắc sau giữ cho WBS vừa đủ chi tiết để giao việc và đo tiến độ:
\begin{itemize}
  \item \textbf{Kích thước:} gói công việc ở cấp thấp nhất làm được trong $\leq$ 3 ngày và có sản phẩm kiểm tra
        được. Quy tắc 8/80 giờ thường dùng được nhóm siết lại thành $\leq$ 24 giờ cho mỗi gói.
  \item \textbf{Trách nhiệm:} mỗi gói có đúng một người chịu trách nhiệm, theo ma trận RACI ở
        Mục~\ref{sec:raci}.
  \item \textbf{Truy vết:} gói công việc ánh xạ 1--1 với task trên Azure Boards; commit trên GitHub ghi mã task,
        ví dụ \texttt{(\#70)}.
\end{itemize}

% ============================================================ Chapter 3
\chapter{Ước lượng effort}
\label{chap:effort}

\textbf{Phương pháp chính là story point với Planning Poker, quy ra giờ bằng velocity thực tế.} Hai phương pháp
đối chiếu, COCOMO~81 cơ bản (tính từ số dòng code thật của repo) và Function Point (tính từ số API và bảng dữ liệu
thật), cho kết quả 3.500--5.100 giờ cho cả sản phẩm nếu làm theo kiểu công ty với một nhóm toàn thời gian. Con số
này dùng để kiểm tra độ hợp lý về quy mô, không thay cho số giờ thực tế của nhóm.

\section{Story point và Planning Poker}

Mỗi buổi Sprint Planning ước lượng theo năm bước:
\begin{enumerate}
  \item PO đọc user story và acceptance criteria trong Product Backlog~\cite{srs}.
  \item Mỗi thành viên chọn một lá bài theo dãy Fibonacci: 1, 2, 3, 5, 8, 13, 21, ?.
  \item Cả nhóm lật bài cùng lúc. Người chọn cao nhất và thấp nhất giải thích lý do; vote lại tối đa 2 vòng.
  \item Story trên 13 điểm phải tách nhỏ. Story mốc tham chiếu là \textbf{\#65 Xuất TXT/SRT = 2 điểm}.
  \item Khi lập kế hoạch, điểm được quy ra giờ theo công thức~\eqref{eq:sp-hours}.
\end{enumerate}
\begin{equation}
  \label{eq:sp-hours}
  \text{Giờ} \approx \text{SP} \times \frac{H_2}{\mathit{SP}_2}
\end{equation}
trong đó $H_2$ và $\mathit{SP}_2$ là tổng giờ thực tế và tổng story point hoàn thành của 2 sprint gần nhất.

\section{Bảng ước lượng theo backlog}

Bảng~\ref{tab:estimate} liệt kê 18 hạng mục theo backlog. Cột ``SP'' là điểm đề xuất để nhóm đối chiếu với điểm
đã vote trên Azure Boards; cột ``Giờ thực tế'' lấy từ board. Giờ ước lượng của từng hạng mục được tính bằng
công thức~\eqref{eq:sp-hours} khi đã có velocity (Mục~\ref{sec:velocity}).

\begin{table}[htbp]
  \centering
  \caption{Ước lượng theo backlog}
  \label{tab:estimate}
  \small
  \begin{tabularx}{\linewidth}{@{}P{2.6cm}LQ{1.6cm}Q{2.6cm}@{}}
    \toprule
    \thead{Task / story} & \thead{Nội dung} & \thead{SP đề xuất} & \thead{Giờ thực tế} \\
    \midrule
    T1.1.1--T1.1.5 & Ghi âm trên trình duyệt, gửi chunk, tạo meeting, test trình duyệt & 13 & \canDien \\
    T1.2.1--T1.2.5 & Ghi trên điện thoại, offline, tự đồng bộ & 13 & \canDien \\
    T1.3.1--T1.3.5 & Upload file lớn, kiểm tra định dạng, chuẩn hoá audio & 8 & \canDien \\
    STT & Whisper, transcript ngôn ngữ gốc, live & 13 & \canDien \\
    \midrule
    \#68 & Model và prompt tóm tắt & 5 & \canDien \\
    \#69 & Transcript dài (map-reduce) & 5 & \canDien \\
    \#70 & Bước tóm tắt, lưu trữ, API & 5 & \canDien \\
    \#71 & Tab Summary và Regenerate & 3 & \canDien \\
    \#61 & Model dịch & 3 & \canDien \\
    \#62 & Lưu bản dịch và API & 3 & \canDien \\
    \#63 & Bước dịch có tiến độ & 3 & \canDien \\
    \#64 & Chọn ngôn ngữ và hiển thị song ngữ & 3 & \canDien \\
    \#65 & Xuất TXT/SRT (mốc tham chiếu) & 2 & \canDien \\
    \#66, \#72 & Đánh giá chất lượng dịch, tóm tắt & 5 & \canDien \\
    \#99, \#100 & Test E2E Sprint 3 & 5 & \canDien \\
    \midrule
    D-1 & App desktop (Electron, đóng gói, bộ cài) & 8 & \canDien \\
    S-1 & Live streaming, STT remote, file dài & 13 & \canDien \\
    O-1 & Deploy AWS (Docker, S3, hướng dẫn) & 8 & \canDien \\
    \midrule
    \textbf{Tổng} & & \textbf{118} & \\
    \bottomrule
  \end{tabularx}
\end{table}

Tổng điểm đề xuất là 118 SP, được cộng lại từ các dòng của bảng; phiên bản 2.0 ghi tổng là 125 SP do cộng sai.
Nhóm tóm tắt và dịch (\#61--\#72) chiếm 37 SP, đúng bằng tổng điểm của hai story \#67 (18 SP) và \#60 (19 SP)
trong Product Backlog~\cite{srs}.

\section{Velocity}
\label{sec:velocity}

Velocity là số story point nhóm hoàn thành trong một sprint. Velocity dùng để lập kế hoạch là trung bình của 3
sprint gần nhất, và từ đó suy ra số sprint còn lại:
\begin{equation}
  \label{eq:velocity}
  v = \frac{1}{3}\sum_{\text{3 sprint gần nhất}} \text{SP hoàn thành},
  \qquad
  \text{Số sprint còn lại} = \frac{\text{SP còn lại trong backlog}}{v}
\end{equation}
Số liệu của ba sprint đầu (Bảng~\ref{tab:velocity}) được lấy từ Azure Boards.

\begin{table}[htbp]
  \centering
  \caption{Velocity theo sprint}
  \label{tab:velocity}
  \begin{tabularx}{\linewidth}{@{}LQ{2.5cm}Q{2.5cm}Q{2.5cm}Q{1.8cm}@{}}
    \toprule
    \thead{Sprint} & \thead{SP cam kết} & \thead{SP hoàn thành} & \thead{Giờ thực tế} & \thead{Giờ/SP} \\
    \midrule
    Sprint 1 & \canDien & \canDien & \canDien & \\
    Sprint 2 & \canDien & \canDien & \canDien & \\
    Sprint 3 & \canDien & \canDien & \canDien & \\
    \midrule
    \textbf{Trung bình} & & & & \\
    \bottomrule
  \end{tabularx}
\end{table}

\section{Đối chiếu bằng COCOMO 81 cơ bản}

Kích thước sản phẩm được đo trong repo ngày 12/11/2026, chỉ tính code viết tay, không tính dòng trống, comment,
test và các component UI sinh tự động (Bảng~\ref{tab:loc}).

\begin{table}[htbp]
  \centering
  \caption{Kích thước mã nguồn ngày 12/11/2026}
  \label{tab:loc}
  \begin{tabular}{@{}P{10cm}r@{}}
    \toprule
    \thead{Thành phần} & \thead{Dòng code} \\
    \midrule
    Backend Python (\path{backend/app}, \path{run.py}) & 3.386 \\
    Frontend TypeScript/React (trừ \path{components/ui}) & 4.518 \\
    Desktop (Electron, build) & 749 \\
    \midrule
    \textbf{Tổng} & \textbf{8.653 (8,65 KLOC)} \\
    \bottomrule
  \end{tabular}
\end{table}

COCOMO~81 cơ bản~\cite{boehm} ước lượng effort $E$ (người-tháng), thời gian $T$ (tháng) và số người $N$ từ kích
thước theo KLOC, với các hệ số $a, b, c, d$ phụ thuộc chế độ dự án:
\begin{equation}
  \label{eq:cocomo}
  E = a \cdot \mathrm{KLOC}^{\,b}, \qquad T = c \cdot E^{\,d}, \qquad N = \frac{E}{T}
\end{equation}
Kết quả cho hai chế độ phù hợp với SumTalk được nêu trong Bảng~\ref{tab:cocomo}, với $E$ tính bằng người-tháng,
$T$ bằng tháng; giờ công quy đổi theo 152 giờ/tháng và chi phí theo 100.000~VND/giờ.

\begin{table}[htbp]
  \centering
  \caption{Ước lượng bằng COCOMO 81 cơ bản (8,65 KLOC)}
  \label{tab:cocomo}
  \small
  \setlength{\tabcolsep}{5pt}
  \begin{tabularx}{\linewidth}{@{}LQ{3.3cm}rrrrr@{}}
    \toprule
    \thead{Chế độ} & \thead{$a$; $b$; $c$; $d$} & \thead{$E$} & \thead{$T$} & \thead{$N$}
      & \thead{Giờ} & \thead{Chi phí (VND)} \\
    \midrule
    Organic & 2,4; 1,05; 2,5; 0,38 & 23,1 & 8,2 & 2,8 & 3.516 & 351,6 triệu \\
    Semi-detached (có AI/ML, nhiều công nghệ) & 3,0; 1,12; 2,5; 0,35 & 33,6 & 8,6 & 3,9 & 5.112 & 511,2 triệu \\
    \bottomrule
  \end{tabularx}
\end{table}

\section{Đối chiếu bằng Function Point}

Function Point được đếm theo IFPUG với trọng số trung bình~\cite{ifpug}, từ hệ thống thật: 27 endpoint API và
6 bảng dữ liệu (Bảng~\ref{tab:fp}).

\begin{table}[htbp]
  \centering
  \caption{Đếm Function Point (IFPUG, trọng số trung bình)}
  \label{tab:fp}
  \small
  \setlength{\tabcolsep}{4pt}
  \begin{tabularx}{\linewidth}{@{}P{2.9cm}LQ{1.3cm}Q{1.4cm}Q{1cm}@{}}
    \toprule
    \thead{Loại} & \thead{Đếm từ hệ thống} & \thead{Số lượng} & \thead{Trọng số} & \thead{Điểm} \\
    \midrule
    EI -- nhập liệu & Tạo meeting, gửi chunk, finalize, đổi tên, xoá, upload (tạo, part, complete, huỷ),
      retranscribe, regenerate/dịch & 11 & 4 & 44 \\
    EO -- xuất có xử lý & Export transcript (5 định dạng), export bản dịch, transcript live, tóm tắt, phát audio,
      API \path{/v1} transcription & 7 & 5 & 35 \\
    EQ -- truy vấn & Danh sách meeting, chi tiết, transcript, health, config, trạng thái upload, danh sách bản dịch,
      bản dịch, models & 9 & 4 & 36 \\
    ILF -- dữ liệu nội bộ & Meeting, Segment, Upload, Summary, Translation, SegmentTranslation & 6 & 10 & 60 \\
    EIF -- dữ liệu ngoài & API LLM, API Whisper/S3 & 2 & 7 & 14 \\
    \midrule
    \textbf{UFP} & & & & \textbf{189} \\
    \bottomrule
  \end{tabularx}
\end{table}

Hệ số điều chỉnh được tính từ tổng điểm của 14 đặc tính hệ thống, giả định là 38 điểm:
\begin{equation*}
  \mathrm{VAF} = 0{,}65 + 0{,}01 \times 38 = 1{,}03, \qquad
  \mathrm{AFP} = 189 \times 1{,}03 \approx 195~\text{FP}
\end{equation*}
Với năng suất 8 giờ/FP, effort khoảng $195 \times 8 \approx$ \textbf{1.560 giờ}.

\section{So sánh và kết luận}

Bảng~\ref{tab:effort-compare} và Hình~\ref{fig:effort-compare} đặt các kết quả cạnh nhau. Ước lượng theo story
point là số dùng để lập kế hoạch, nhưng còn chờ giờ thực tế để quy đổi.

\begin{table}[htbp]
  \centering
  \caption{So sánh các phương pháp ước lượng}
  \label{tab:effort-compare}
  \begin{tabularx}{\linewidth}{@{}P{5.6cm}P{3.4cm}L@{}}
    \toprule
    \thead{Phương pháp} & \thead{Effort} & \thead{Ghi chú} \\
    \midrule
    Story point (118 SP $\times$ giờ/SP thực tế) & \canDien & Số dùng để lập kế hoạch \\
    Function Point (195 FP $\times$ 8 giờ) & Khoảng 1.560 giờ & Chỉ phần chức năng, chưa gồm desktop và deploy \\
    COCOMO organic & Khoảng 3.500 giờ & Mô hình công ty, gồm cả quản lý và tài liệu \\
    COCOMO semi-detached & Khoảng 5.100 giờ & Như trên, hệ số cho dự án nhiều công nghệ \\
    \bottomrule
  \end{tabularx}
\end{table}

\begin{figure}[htbp]
  \centering
  \begin{tikzpicture}
    \begin{axis}[
        width=0.72\linewidth, height=4.6cm,
        xbar, bar width=0.5cm,
        xmin=0, xmax=6400,
        symbolic y coords={FP,CO,CS}, ytick=data, y dir=reverse,
        yticklabels={{Function Point},{COCOMO organic},{COCOMO semi-detached}},
        xlabel={Giờ công ước lượng},
        nodes near coords, every node near coord/.append style={font=\footnotesize},
        grid=major, ymajorgrids=false, grid style={gray!25},
        tick label style={font=\small}, enlarge y limits=0.3,
        /pgf/number format/.cd, use comma, 1000 sep={.}]
      \addplot[fill=primaryblue, draw=primaryblue] coordinates {(1560,FP) (3516,CO) (5112,CS)};
    \end{axis}
  \end{tikzpicture}
  \caption{Effort theo hai phương pháp đối chiếu}
  \label{fig:effort-compare}
\end{figure}

Giờ thực tế của nhóm thấp hơn nhiều so với COCOMO, vì bốn lý do: nhóm dùng framework và thư viện có sẵn (FastAPI,
React, faster-whisper, Electron); có công cụ AI hỗ trợ viết code; phạm vi POC chưa có tài khoản và phân quyền; và
COCOMO~81 được hiệu chỉnh trên các dự án của những năm 1970--80. Vì vậy, khi báo cáo, story point là số kế hoạch;
Function Point và COCOMO chỉ dùng để chứng minh quy mô sản phẩm.

% ============================================================ Chapter 4
\chapter{Lịch trình}
\label{chap:schedule}

\textbf{Đường găng là A $\rightarrow$ B $\rightarrow$ C/F $\rightarrow$ G $\rightarrow$ U $\rightarrow$ V
$\rightarrow$ W, dài 33 ngày làm việc} (khoảng 7 tuần nếu đủ người làm song song). Phần tóm tắt và dịch (H--S) có
3--12 ngày dự trữ, nên một thành viên chuyển sang hỗ trợ phần live và STT khi cần sẽ không làm trễ dự án.

Thời lượng của từng công việc là \textbf{ước lượng} từ story point (Chương~\ref{chap:effort}). Phụ thuộc lấy từ
board (ví dụ \#61 $\rightarrow$ \#63, \#68 $\rightarrow$ \#70) và từ kiến trúc. Ngày bắt đầu của Gantt giả định là
07/09/2026: \canDien[thay bằng ngày bắt đầu Sprint 1 thật]. Toàn bộ bảng CPM được tính lại bằng
\texttt{pm\_calc.py}~\cite{pmcalc}.

\section{Danh sách công việc và phương pháp đường găng}

Với mỗi công việc có thời lượng $d$, phương pháp đường găng tính thời điểm bắt đầu và kết thúc sớm nhất (ES, EF)
bằng một lượt đi xuôi, thời điểm muộn nhất (LS, LF) bằng một lượt đi ngược, rồi lấy hiệu làm dự trữ:
\begin{equation}
  \label{eq:cpm}
  \begin{aligned}
    \mathrm{ES} &= \max_{\text{trước}} \mathrm{EF}, & \mathrm{EF} &= \mathrm{ES} + d, \\
    \mathrm{LF} &= \min_{\text{sau}} \mathrm{LS}, & \mathrm{LS} &= \mathrm{LF} - d, \qquad
    \text{Dự trữ} = \mathrm{LS} - \mathrm{ES}.
  \end{aligned}
\end{equation}
Công việc có dự trữ bằng 0 là công việc găng: trễ bao nhiêu ngày thì cả dự án trễ bấy nhiêu ngày. Kết quả cho 23
công việc được trình bày trong Bảng~\ref{tab:cpm}.

\begin{table}[htbp]
  \centering
  \caption{Danh sách công việc, phụ thuộc và kết quả CPM (ngày làm việc)}
  \label{tab:cpm}
  \small
  \setlength{\tabcolsep}{4pt}
  \begin{tabularx}{\linewidth}{@{}Q{0.6cm}LQ{1.0cm}P{2.1cm}rrrrQ{1.2cm}@{}}
    \toprule
    \thead{ID} & \thead{Công việc} & \thead{Ngày} & \thead{Phụ thuộc} & \thead{ES} & \thead{EF} & \thead{LS}
      & \thead{LF} & \thead{Dự trữ} \\
    \midrule
    \textbf{A} & Khởi động: charter, backlog, yêu cầu và UI/UX & 5 & -- & 0 & 5 & 0 & 5 & \textbf{0} \\
    \textbf{B} & Kiến trúc và môi trường dev & 3 & A & 5 & 8 & 5 & 8 & \textbf{0} \\
    \textbf{C} & T1.1 Ghi âm trên trình duyệt, gửi chunk & 8 & B & 8 & 16 & 8 & 16 & \textbf{0} \\
    D & T1.2 Offline/mobile (IndexedDB, đồng bộ) & 6 & C & 16 & 22 & 24 & 30 & 8 \\
    E & T1.3 Upload file lớn, chuẩn hoá audio & 6 & B & 8 & 14 & 24 & 30 & 16 \\
    \textbf{F} & STT Whisper, transcript ngôn ngữ gốc & 8 & B & 8 & 16 & 8 & 16 & \textbf{0} \\
    \textbf{G} & Phiên âm live & 5 & C, F & 16 & 21 & 16 & 21 & \textbf{0} \\
    H & Benchmark model LLM (Sprint 2) & 4 & A & 5 & 9 & 10 & 14 & 5 \\
    I & \#68 Model và prompt tóm tắt & 3 & H & 9 & 12 & 16 & 19 & 7 \\
    J & \#69 Transcript dài (map-reduce) & 3 & I & 12 & 15 & 24 & 27 & 12 \\
    K & \#70 Bước tóm tắt, lưu trữ, API & 4 & I, F & 16 & 20 & 19 & 23 & 3 \\
    L & \#71 Tab Summary và Regenerate & 3 & K & 20 & 23 & 23 & 26 & 3 \\
    M & \#61 Model dịch & 3 & H & 9 & 12 & 14 & 17 & 5 \\
    N & \#62 Lưu bản dịch và API & 3 & M & 12 & 15 & 17 & 20 & 5 \\
    O & \#63 Bước dịch có tiến độ & 3 & M, N & 15 & 18 & 20 & 23 & 5 \\
    P & \#64 Chọn ngôn ngữ và song ngữ & 3 & O & 18 & 21 & 23 & 26 & 5 \\
    Q & \#65 Xuất TXT/SRT & 2 & N & 15 & 17 & 24 & 26 & 9 \\
    R & \#66/\#72 Đánh giá chất lượng & 3 & J, O & 18 & 21 & 27 & 30 & 9 \\
    S & \#99/\#100 Test E2E Sprint 3 & 4 & L, P, Q & 23 & 27 & 26 & 30 & 3 \\
    T & App desktop (Electron, PyInstaller) & 5 & G & 21 & 26 & 25 & 30 & 4 \\
    \textbf{U} & Streaming, STT remote, file dài & 5 & G & 21 & 26 & 21 & 26 & \textbf{0} \\
    \textbf{V} & Deploy AWS (Docker, S3, hướng dẫn) & 4 & U & 26 & 30 & 26 & 30 & \textbf{0} \\
    \textbf{W} & UAT, phát hành, đóng dự án & 3 & S, T, V, R, D, E & 30 & 33 & 30 & 33 & \textbf{0} \\
    \bottomrule
  \end{tabularx}
\end{table}

Tám công việc găng được in đậm. Đường găng có hai nhánh song song cùng dài 8 ngày ở giữa (C ghi âm và F STT), nên
trễ ở bất kỳ nhánh nào cũng làm trễ G. Ngược lại, E (upload) có tới 16 ngày dự trữ và J (map-reduce) có 12 ngày:
đây là nơi có thể tạm lấy người sang hỗ trợ công việc găng.

\section{Sơ đồ mạng}

Hình~\ref{fig:network} vẽ 23 công việc dưới dạng sơ đồ mạng AON: mỗi nút là một công việc, mũi tên là quan hệ
phụ thuộc kết thúc -- bắt đầu. Đường găng được tô đỏ. Sơ đồ cho thấy dự án có hai luồng gần như độc lập sau bước
khởi động: luồng ghi âm, STT, live và triển khai ở nửa trên (chứa toàn bộ đường găng), và luồng LLM cho tóm tắt,
dịch ở nửa dưới, chỉ gặp luồng trên tại K (\#70 cần STT) và tại W.

\begin{figure}[htbp]
  \centering
  \begin{tikzpicture}[
      nd/.style={draw=primaryblue,fill=accentblue,rounded corners=1.5pt,align=center,font=\scriptsize,
                 text width=1.45cm,minimum height=0.9cm,inner sep=2pt,
                 execute at begin node=\setstretch{1}},
      cr/.style={nd,draw=todored,fill=todored!10,line width=0.9pt},
      ar/.style={-{Stealth[length=1.8mm]},primaryblue},
      ca/.style={-{Stealth[length=2mm]},todored,line width=1pt}]
    % Nodes: upper half = recording/STT/deploy stream, lower half = LLM stream.
    \node[cr] (A) at (0,0)     {\textbf{A} (5)\\ Khởi động};
    \node[cr] (B) at (2,1.4)   {\textbf{B} (3)\\ Kiến trúc};
    \node[nd] (E) at (4,4.2)   {\textbf{E} (6)\\ Upload};
    \node[cr] (C) at (4,2.8)   {\textbf{C} (8)\\ Ghi âm};
    \node[nd] (D) at (6,2.8)   {\textbf{D} (6)\\ Offline};
    \node[cr] (F) at (4,1.4)   {\textbf{F} (8)\\ STT gốc};
    \node[cr] (G) at (6,1.4)   {\textbf{G} (5)\\ Live};
    \node[cr] (U) at (8,1.4)   {\textbf{U} (5)\\ Streaming};
    \node[cr] (V) at (10,1.4)  {\textbf{V} (4)\\ Deploy AWS};
    \node[nd] (T) at (8,0)     {\textbf{T} (5)\\ Desktop};
    \node[nd] (H) at (2,-1.4)  {\textbf{H} (4)\\ Benchmark};
    \node[nd] (I) at (4,-1.4)  {\textbf{I} (3)\\ \#68};
    \node[nd] (K) at (6,-1.4)  {\textbf{K} (4)\\ \#70};
    \node[nd] (L) at (8,-1.4)  {\textbf{L} (3)\\ \#71};
    \node[nd] (S) at (12,-1.4) {\textbf{S} (4)\\ \#99/\#100};
    \node[nd] (J) at (6,-2.8)  {\textbf{J} (3)\\ \#69};
    \node[nd] (R) at (10,-2.8) {\textbf{R} (3)\\ \#66/\#72};
    \node[nd] (M) at (4,-4.2)  {\textbf{M} (3)\\ \#61};
    \node[nd] (O) at (8,-4.2)  {\textbf{O} (3)\\ \#63};
    \node[nd] (P) at (10,-4.2) {\textbf{P} (3)\\ \#64};
    \node[nd] (N) at (6,-5.6)  {\textbf{N} (3)\\ \#62};
    \node[nd] (Q) at (8,-5.6)  {\textbf{Q} (2)\\ \#65};
    \node[cr,minimum height=7.9cm] (W) at (14,0.7) {\textbf{W} (3)\\ UAT, phát hành};
    % Critical path
    \draw[ca] (A) -- (B);
    \draw[ca] (B) -- (C);
    \draw[ca] (B) -- (F);
    \draw[ca] (C) -- (G);
    \draw[ca] (F) -- (G);
    \draw[ca] (G) -- (U);
    \draw[ca] (U) -- (V);
    \draw[ca] (V) -- (V -| W.west);
    % Other dependencies
    \draw[ar] (A) -- (H);
    \draw[ar] (B) |- (E);
    \draw[ar] (C) -- (D);
    \draw[ar] (F) -- (K);
    \draw[ar] (G) -- (T);
    \draw[ar] (E) -- (E -| W.west);
    \draw[ar] (D) -- (D -| W.west);
    \draw[ar] (T) -- (T -| W.west);
    \draw[ar] (S) -- (S -| W.west);
    \draw[ar] (H) -- (I);
    \draw[ar] (H) |- (M);
    \draw[ar] (I) -- (K);
    \draw[ar] (I) -- (J);
    \draw[ar] (K) -- (L);
    \draw[ar] (L) -- (S);
    \draw[ar] (J) -- (R);
    \draw[ar] (M) -- (O);
    \draw[ar] (M) -- (N);
    \draw[ar] (N) -- (O);
    \draw[ar] (N) -- (Q);
    \draw[ar] (O) -- (P);
    \draw[ar] (O) -- (R);
    \draw[ar] (R) -- (R -| W.west);
    % P and Q join S from below; the line bridges over R's arrow.
    \draw[white,line width=3pt] (12,-2.62) -- (12,-2.98);
    \draw[ar] (Q.east) -| (S.south);
    \draw[primaryblue] (P.east) -- (P.east -| S.south);
    \fill[primaryblue] (P.east -| S.south) circle (1.3pt);
  \end{tikzpicture}
  \caption{Sơ đồ mạng AON; số trong ngoặc là thời lượng (ngày làm việc), đường găng tô đỏ}
  \label{fig:network}
\end{figure}

\section{Biểu đồ Gantt}

Biểu đồ Gantt (Hình~\ref{fig:gantt}) đặt các công việc lên trục thời gian theo ES và EF của Bảng~\ref{tab:cpm},
tính bằng ngày làm việc (bỏ cuối tuần) kể từ ngày bắt đầu giả định 07/09/2026. Phần nét đứt sau mỗi thanh không găng
là dự trữ: công việc có thể kết thúc muộn tới đó mà không làm trễ W.

\begin{figure}[htbp]
  \centering
  \newcommand{\gcrit}[3]{\ganttbar[bar/.append style={fill=todored!80,draw=todored}]{#1}{#2}{#3}}
  \newcommand{\gtask}[3]{\ganttbar[bar/.append style={fill=primaryblue!70,draw=primaryblue}]{#1}{#2}{#3}}
  \newcommand{\gslack}[2]{\ganttbar[bar/.append style={fill=gray!12,draw=chartgray,dashed},bar height=0.3]{}{#1}{#2}}
  \begin{ganttchart}[
      x unit=0.3cm, y unit chart=0.47cm, y unit title=0.5cm,
      hgrid={draw=gray!15}, vgrid={*4{draw=none},*1{gray!35}},
      title/.append style={fill=accentblue,draw=primaryblue!50},
      title label font=\scriptsize, bar label font=\scriptsize,
      bar height=0.55, bar top shift=0.225,
      bar label node/.append style={align=right}]
                {1}{35}
    \gantttitle{Tuần 1}{5} \gantttitle{Tuần 2}{5} \gantttitle{Tuần 3}{5} \gantttitle{Tuần 4}{5}
    \gantttitle{Tuần 5}{5} \gantttitle{Tuần 6}{5} \gantttitle{Tuần 7}{5} \\
    \gantttitle{07/09}{5} \gantttitle{14/09}{5} \gantttitle{21/09}{5} \gantttitle{28/09}{5}
    \gantttitle{05/10}{5} \gantttitle{12/10}{5} \gantttitle{19/10}{5} \\
    \gcrit{A -- Charter, backlog, yêu cầu}{1}{5} \\
    \gcrit{B -- Kiến trúc, môi trường}{6}{8} \\
    \gcrit{C -- Ghi âm trình duyệt}{9}{16} \\
    \gtask{E -- Upload file lớn}{9}{14} \gslack{15}{30} \\
    \gcrit{F -- STT ngôn ngữ gốc}{9}{16} \\
    \gtask{D -- Offline/mobile}{17}{22} \gslack{23}{30} \\
    \gcrit{G -- Phiên âm live}{17}{21} \\
    \gtask{H -- Benchmark LLM}{6}{9} \gslack{10}{14} \\
    \gtask{I -- \#68 Model, prompt}{10}{12} \gslack{13}{19} \\
    \gtask{J -- \#69 Map-reduce}{13}{15} \gslack{16}{27} \\
    \gtask{K -- \#70 Bước tóm tắt, API}{17}{20} \gslack{21}{23} \\
    \gtask{L -- \#71 Tab Summary}{21}{23} \gslack{24}{26} \\
    \gtask{M -- \#61 Model dịch}{10}{12} \gslack{13}{17} \\
    \gtask{N -- \#62 Lưu bản dịch, API}{13}{15} \gslack{16}{20} \\
    \gtask{O -- \#63 Bước dịch}{16}{18} \gslack{19}{23} \\
    \gtask{P -- \#64 Song ngữ}{19}{21} \gslack{22}{26} \\
    \gtask{Q -- \#65 Xuất TXT/SRT}{16}{17} \gslack{18}{26} \\
    \gtask{R -- \#66/\#72 Đánh giá}{19}{21} \gslack{22}{30} \\
    \gtask{S -- \#99/\#100 Test E2E}{24}{27} \gslack{28}{30} \\
    \gtask{T -- App desktop}{22}{26} \gslack{27}{30} \\
    \gcrit{U -- Streaming, STT remote}{22}{26} \\
    \gcrit{V -- Deploy AWS}{27}{30} \\
    \gcrit{W -- UAT, phát hành}{31}{33}
  \end{ganttchart}
  \caption{Biểu đồ Gantt theo CPM (ngày bắt đầu 07/09/2026 là giả định); đỏ: công việc găng, nét đứt: dự trữ}
  \label{fig:gantt}
\end{figure}

\section{Release plan theo sprint}

Các công việc được gom thành bốn sprint và một giai đoạn kết thúc, mỗi giai đoạn ra một bản phát hành
(Bảng~\ref{tab:release}). Story \#60 được dời vào Sprint 3 theo Change Log~\cite{mcr}.

\begin{table}[htbp]
  \centering
  \caption{Release plan theo sprint}
  \label{tab:release}
  \small
  \begin{tabularx}{\linewidth}{@{}P{1.9cm}LP{3.6cm}Q{1.6cm}@{}}
    \toprule
    \thead{Sprint} & \thead{Mục tiêu sprint} & \thead{Công việc} & \thead{Phát hành} \\
    \midrule
    Sprint 1 & Ghi âm và transcript gốc chạy được (POC E1) & A, B, C, D, E, F & v0.1 \\
    Sprint 2 & Transcript live; chọn model LLM & G, H & v0.2 \\
    Sprint 3 & Tóm tắt và dịch hoàn chỉnh, có test E2E & I--S (gồm \#60 được dời vào) & v0.3 \\
    Sprint 4 & App desktop, live dạng streaming, STT qua server, bản web trên AWS & T, U, V & v0.4 \\
    Kết thúc & UAT, nghiệm thu & W & v1.0 \\
    \bottomrule
  \end{tabularx}
\end{table}

Cột ``Công việc'' là kế hoạch; \canDien[đối chiếu với việc đã thực sự làm trong từng sprint trên Azure Boards].

\section{Cách dùng lịch trình khi theo dõi}

\begin{itemize}
  \item \textbf{Công việc găng trễ:} trễ 1 ngày thì cả dự án trễ 1 ngày. Người phụ trách báo ngay trong Daily; nếu
        phải dời mốc thì ghi vào Change Log~\cite{mcr}.
  \item \textbf{Dùng dự trữ:} công việc có dự trữ lớn (E: 16 ngày, J: 12 ngày) là nơi có thể tạm lấy người sang
        hỗ trợ công việc găng.
  \item \textbf{Cập nhật:} mỗi sprint ghi ngày thực tế vào bảng, rồi tính lại bằng lệnh
        \texttt{python docs/project/tools/pm\_calc.py}.
\end{itemize}

% ============================================================ Chapter 5
\chapter{Quản lý rủi ro}
\label{chap:risk}

\textbf{Rủi ro lớn nhất hiện tại là chi phí và giới hạn của API bên ngoài khi deploy (R6, R8), và độ chính xác STT
với tiếng Việt hoặc họp trộn ngôn ngữ (R1).} Cả hai đã có biện pháp kỹ thuật nhưng cần được đo lại trong đợt UAT.
Risk Register hiện có 14 rủi ro, 5 rủi ro ở vùng Cao.

\section{Thang đánh giá}

Mỗi rủi ro được chấm xác suất (P) và tác động (I) theo thang 1--5 (Bảng~\ref{tab:risk-scale}).

% [H]: a longtable follows right after this section; a floating table would leave half of this page empty.
\begin{table}[H]
  \centering
  \caption{Thang điểm xác suất và tác động}
  \label{tab:risk-scale}
  \begin{tabularx}{\linewidth}{@{}Q{1.2cm}P{4.4cm}L@{}}
    \toprule
    \thead{Điểm} & \thead{Xác suất (P)} & \thead{Tác động (I)} \\
    \midrule
    1 & Hiếm ($<$ 10\%) & Không đáng kể, xử lý trong ngày \\
    2 & Thấp (10--30\%) & Trễ $<$ 2 ngày, không ảnh hưởng mốc \\
    3 & Trung bình (30--50\%) & Trễ 1 sprint một phần việc, hoặc tốn thêm $<$ 10\% ngân sách \\
    4 & Cao (50--80\%) & Lỡ mốc, phải cắt phạm vi, hoặc tốn thêm 10--20\% ngân sách \\
    5 & Gần như chắc chắn ($>$ 80\%) & Mục tiêu dự án thất bại \\
    \bottomrule
  \end{tabularx}
\end{table}

Mức độ của rủi ro là tích hai điểm và quyết định cách xử lý:
\begin{equation}
  \label{eq:risk-level}
  \text{Mức độ} = P \times I, \qquad
  \begin{cases}
    1\text{--}5: & \text{Thấp -- theo dõi,} \\
    6\text{--}11: & \text{Trung bình -- có kế hoạch ứng phó,} \\
    12\text{--}25: & \text{Cao -- xử lý ngay, báo Sponsor.}
  \end{cases}
\end{equation}

\section{Risk Register}

Risk Register được chia thành hai bảng có cùng mã rủi ro. Bảng~\ref{tab:risk-assess} ghi mô tả, điểm đánh giá
và trạng thái; Bảng~\ref{tab:risk-response} ghi chiến lược, biện pháp ứng phó, dấu hiệu cảnh báo và người phụ
trách.

R1--R6 là các rủi ro đã nêu trong buổi họp của nhóm. Repo chỉ ghi rõ R6 (ngân sách API/LLM); nội dung R1--R5 trong
bảng là đề xuất: \canDien[đối chiếu với biên bản và sửa lại cho khớp tên, mô tả gốc]. R7 trở đi là các rủi ro mới
phát hiện khi làm app desktop, streaming và deploy.

\begingroup\small
\setlength{\tabcolsep}{4pt}
\begin{longtable}{@{}P{0.9cm}P{6.3cm}P{1.9cm}Q{0.5cm}Q{0.5cm}Q{0.9cm}P{2.7cm}@{}}
  \caption{Risk Register: đánh giá và trạng thái}\label{tab:risk-assess}\\
  \toprule
  \thead{ID} & \thead{Rủi ro} & \thead{Hạng mục} & \thead{P} & \thead{I} & \thead{Mức} & \thead{Trạng thái} \\
  \midrule
  \endfirsthead
  \multicolumn{7}{@{}l}{\footnotesize\itshape Bảng~\thetable{} (tiếp theo)}\\
  \toprule
  \thead{ID} & \thead{Rủi ro} & \thead{Hạng mục} & \thead{P} & \thead{I} & \thead{Mức} & \thead{Trạng thái} \\
  \midrule
  \endhead
  \bottomrule
  \endlastfoot
  R1 & STT sai với tiếng Việt, họp trộn Anh -- Việt, hoặc tự sinh chữ khi im lặng & Kỹ thuật & 3 & 4 & \textbf{12}
     & Đang theo dõi \\
  R2 & Máy người dùng không có GPU, CPU quá chậm cho live & Kỹ thuật & 4 & 3 & \textbf{12}
     & \textbf{Đã xử lý} (Sprint 4) \\
  R3 & Trễ tiến độ do làm bán thời gian, trùng lịch công việc khác & Lịch & 4 & 3 & \textbf{12} & Đang theo dõi \\
  R4 & Mất dữ liệu ghi âm khi rớt mạng, đóng tab hoặc app & Kỹ thuật & 2 & 4 & 8 & Đã xử lý \\
  R5 & Yêu cầu thay đổi, phình phạm vi (thêm desktop, streaming, cloud) & Phạm vi & 4 & 3 & \textbf{12}
     & Đang theo dõi \\
  R6 & Chi phí API LLM/STT vượt ngân sách & Chi phí & 3 & 3 & 9 & Đang theo dõi \\
  R7 & Model LLM chạy local (Ollama) quá nặng cho máy nhóm & Kỹ thuật & 3 & 2 & 6 & Đã xử lý \\
  R8 & Rate limit của Whisper API (lỗi 429) khi ghi live & Bên ngoài & 3 & 3 & 9 & Đang theo dõi \\
  R9 & Lộ dữ liệu họp khi deploy công khai (POC chưa có tài khoản) & Bảo mật & 2 & 5 & 10
     & Đã xử lý (bản test) \\
  R10 & Chi phí AWS phát sinh do quên tắt máy GPU hoặc Elastic IP & Chi phí & 3 & 3 & 9 & Mới \\
  R11 & Mất quyền truy cập repo/CI (push bị 403), mất code chưa push & Cấu hình & 3 & 4 & \textbf{12}
      & \textbf{Đang xảy ra} (12/11/2026) \\
  R12 & Bộ cài Windows chưa ký số, bị SmartScreen chặn & Triển khai & 4 & 2 & 8 & Chấp nhận \\
  R13 & iOS dừng micro khi khoá màn hình & Kỹ thuật & 5 & 2 & 10 & Chấp nhận \\
  R14 & Thành viên chủ chốt vắng hoặc rời nhóm & Nhân sự & 2 & 4 & 8 & Đang theo dõi \\
\end{longtable}
\endgroup

\begingroup\small
\setlength{\tabcolsep}{4pt}
\begin{longtable}{@{}P{0.9cm}P{1.9cm}P{6.4cm}P{3.0cm}P{2.1cm}@{}}
  \caption{Risk Register: ứng phó và theo dõi}\label{tab:risk-response}\\
  \toprule
  \thead{ID} & \thead{Chiến lược} & \thead{Biện pháp ứng phó} & \thead{Dấu hiệu cảnh báo} & \thead{Phụ trách} \\
  \midrule
  \endfirsthead
  \multicolumn{5}{@{}l}{\footnotesize\itshape Bảng~\thetable{} (tiếp theo)}\\
  \toprule
  \thead{ID} & \thead{Chiến lược} & \thead{Biện pháp ứng phó} & \thead{Dấu hiệu cảnh báo} & \thead{Phụ trách} \\
  \midrule
  \endhead
  \bottomrule
  \endlastfoot
  R1 & Giảm & Nhận diện ngôn ngữ theo từng câu; lọc chữ ``ảo''; cho chọn cố định ngôn ngữ và phiên âm lại; dùng
       \texttt{large-v3-turbo} qua server GPU & WER cao trong UAT, người dùng sửa nhiều & \role{Dev STT} \\
  R2 & Giảm & Engine remote (\texttt{STT\_PROVIDER=remote}): server GPU hoặc Whisper API, client không tải model;
       bộ cài ghi sẵn cấu hình & Live trễ $>$ 5\,s & \role{Dev} \\
  R3 & Giảm & Sprint ngắn; ưu tiên việc găng (Chương~\ref{chap:schedule}); dùng dự trữ của nhánh tóm tắt và dịch
     & Burndown nằm trên đường lý tưởng 2 ngày liền & PM \\
  R4 & Tránh & Lưu chunk vào IndexedDB trước khi gửi, tự đồng bộ; khôi phục sau crash; hỏi trước khi đóng khi đang
       ghi & Lỗi đồng bộ trong log & \role{Dev frontend} \\
  R5 & Chấp nhận có kiểm soát & Mọi thay đổi qua PO và Change Log~\cite{mcr}; việc ngoài phạm vi để sau v1.0
     & Story mới chen vào giữa sprint & PO \\
  R6 & Giảm & Log token và chi phí từng lần chạy; \texttt{AI\_MOCK} khi demo; \texttt{LIVE\_INTERVAL\_SEC}; VAD bỏ
       đoạn im lặng không gửi; AWS Budget cảnh báo & Chi phí tuần $>$ 25\% ngân sách tháng & PM \\
  R7 & Chuyển giao & Client chuẩn OpenAI: đổi sang LLM cloud bằng cấu hình & Tóm tắt $>$ 2 phút hoặc hết RAM
     & \role{Dev AI} \\
  R8 & Giảm & Retry có back-off; \texttt{LIVE\_INTERVAL\_SEC}; giới hạn \texttt{STT\_CONCURRENCY}; chuyển sang server
       GPU riêng & Lỗi 429 trong log & \role{Dev} \\
  R9 & Giảm & HTTPS và đăng nhập (Caddy); bucket S3 private, presigned URL ngắn hạn; IAM role tối thiểu;
       \path{/v1} có API key; không commit \texttt{.env} & Truy cập lạ trong log Caddy & \role{Dev ops} \\
  R10 & Giảm & AWS Budget gửi email; tắt máy khi không test; checklist gỡ bỏ trong hướng dẫn deploy
      & Hoá đơn tăng bất thường & PM \\
  R11 & Giảm & Thêm collaborator đúng tài khoản; push sau mỗi task; nhánh feature theo sprint
      & \texttt{git push} lỗi & PM \\
  R12 & Chấp nhận & Hướng dẫn ``More info $\rightarrow$ Run anyway''; demo bằng bản web; ký số khi thương mại
      & Người dùng báo không cài được & \role{Dev desktop} \\
  R13 & Chấp nhận & Giữ màn hình sáng (Wake Lock); ghi chú hạn chế; app native ở phiên bản sau & -- & PO \\
  R14 & Giảm & Code review chéo; tài liệu README/ARCHITECTURE; HANDOFF khi bàn giao & Task bị kẹt $>$ 2 ngày & PM \\
\end{longtable}
\endgroup

Trong 14 rủi ro, chín rủi ro được xử lý bằng chiến lược \emph{giảm}, ba rủi ro được \emph{chấp nhận} (R5 có kiểm
soát, R12 và R13 là hạn chế đã biết), một rủi ro được \emph{tránh} (R4) và một được \emph{chuyển giao} (R7). R11 là
rủi ro duy nhất đang xảy ra tại thời điểm lập kế hoạch.

\section{Ma trận rủi ro}

Hình~\ref{fig:risk-matrix} đặt 14 rủi ro lên ma trận xác suất $\times$ tác động theo điểm \emph{trước} ứng phó.
Năm rủi ro ở vùng Cao ($\geq$ 12) là R1, R2, R3, R5 và R11. Sau khi có engine remote, R2 đã giảm xuống mức
$2 \times 2 = 4$ (mũi tên nét đứt).

\begin{figure}[htbp]
  \centering
  \begin{tikzpicture}[
      cell/.style={draw=white,line width=1.2pt,minimum width=2.3cm,minimum height=1.15cm,
                   align=center,font=\footnotesize,inner sep=1pt}]
    \foreach \p in {1,...,5} {
      \foreach \i in {1,...,5} {
        \pgfmathtruncatemacro{\lv}{\p*\i}
        \ifnum\lv>11 \def\cellcol{todored!25}\else\ifnum\lv>5 \def\cellcol{accentblue}\else\def\cellcol{gray!10}\fi\fi
        \node[cell,fill=\cellcol] (c\p\i) at ({(\i-1)*2.3},{(\p-1)*1.15}) {};
        \node[font=\tiny,text=chartgray,anchor=north east,inner sep=2pt] at (c\p\i.north east) {\lv};
      }
    }
    \begin{scope}[every node/.style={font=\footnotesize}]
      \node at (c52) {R13};
      \node at (c42) {R12};
      \node at (c43) {R2, R3, R5};
      \node at (c32) {R7};
      \node at (c33) {R6, R8, R10};
      \node at (c34) {\textbf{R1, R11}};
      \node at (c24) {R4, R14};
      \node at (c25) {R9};
      \node[text=chartgray,font=\footnotesize\itshape] at (c22) {R2 (sau)};
    \end{scope}
    \draw[-{Stealth[length=2mm]},chartgray,dashed,thick] ($(c43.south west)+(0.3,0.2)$) -- ($(c22.north east)+(-0.4,-0.3)$);
    % Axes
    \foreach \p/\lab in {1/Hiếm,2/Thấp,3/Trung bình,4/Cao,5/Gần chắc chắn}
      \node[anchor=east,font=\footnotesize,align=right] at ($(c\p1.west)+(-0.1,0)$) {\p\ -- \lab};
    \foreach \i in {1,...,5}
      \node[anchor=north,font=\footnotesize] at ($(c1\i.south)+(0,-0.05)$) {\i};
    \node[font=\small\bfseries,anchor=north] at ($(c13.south)+(0,-0.5)$) {Tác động (I)};
    \node[font=\small\bfseries,rotate=90,anchor=south] at ($(c31.west)+(-3.1,0)$) {Xác suất (P)};
    % Legend
    \foreach \col/\lab [count=\k] in {gray!10/Thấp (1--5),accentblue/Trung bình (6--11),todored!25/Cao (12--25)}
      \node[draw=gray!50,fill=\col,minimum width=0.45cm,minimum height=0.3cm,
            label={[font=\footnotesize]right:\lab}] at ({0.6+(\k-1)*3.9},-1.95) {};
  \end{tikzpicture}
  \caption{Ma trận rủi ro theo xác suất và tác động, trước ứng phó}
  \label{fig:risk-matrix}
\end{figure}

\section{Quy trình quản lý rủi ro}

Rủi ro được quản lý liên tục theo năm bước, gắn với các sự kiện Scrum:
\begin{enumerate}
  \item \textbf{Nhận diện:} ở mỗi Sprint Planning và Retrospective, nhóm trả lời câu hỏi ``Điều gì có thể làm
        sprint sau thất bại?''.
  \item \textbf{Đánh giá:} chấm P và I theo Bảng~\ref{tab:risk-scale}, cập nhật Risk Register.
  \item \textbf{Ứng phó:} chọn một trong bốn chiến lược Tránh, Giảm, Chuyển giao, Chấp nhận; ghi biện pháp và
        người phụ trách.
  \item \textbf{Theo dõi:} người phụ trách báo tình trạng trong Daily khi dấu hiệu cảnh báo xuất hiện; PM rà toàn
        bộ bảng ở mỗi Sprint Review.
  \item \textbf{Đóng:} rủi ro đã qua hoặc đã xử lý xong được chuyển sang ``Đã đóng'', bài học được ghi vào Lessons
        Learned~\cite{pcr}.
\end{enumerate}

\section{Dự phòng ngân sách và thời gian}

Dự án giữ hai loại dự phòng cho các rủi ro đã nhận diện. Về ngân sách, 10\% tổng ngân sách
(Chương~\ref{chap:cost}) được dành cho các rủi ro chi phí R6, R8 và R10. Về thời gian, dự trữ 3--12 ngày của
nhánh tóm tắt và dịch cùng 3 ngày UAT cuối dự án là vùng đệm cho trễ tiến độ.

% ============================================================ Chapter 6
\chapter{Nguồn lực}
\label{chap:resource}

\textbf{Ràng buộc chính là giờ rảnh của từng người, không phải số người.} Nguồn lực của dự án gồm con người (tính
theo giờ mỗi sprint), máy móc (laptop GPU của nhóm, máy chủ AWS) và dịch vụ (API Whisper và LLM). Việc san tải dựa
vào dự trữ của các công việc không găng ở Chương~\ref{chap:schedule}.

\section{Nguồn lực con người theo sprint}

Năng lực của mỗi người trong một sprint được tính theo công thức~\eqref{eq:capacity}, trong đó hệ số 0,8 trừ thời
gian họp và gián đoạn:
\begin{equation}
  \label{eq:capacity}
  \text{Năng lực (giờ)} = \text{số ngày của sprint} \times \text{giờ rảnh mỗi ngày} \times 0{,}8
\end{equation}
Bảng~\ref{tab:resource-hours} là khung ghi giờ theo người và sprint;
\canDien[lấy từ số liệu tải theo người trên Azure Boards hoặc dashboard chi phí].

\begin{table}[htbp]
  \centering
  \caption{Giờ làm việc theo thành viên và sprint}
  \label{tab:resource-hours}
  \small
  \begin{tabularx}{\linewidth}{@{}P{2.7cm}LQ{1.4cm}Q{1.4cm}Q{1.4cm}Q{1.4cm}Q{1.2cm}@{}}
    \toprule
    \thead{Thành viên} & \thead{Vai trò} & \thead{Sprint 1} & \thead{Sprint 2} & \thead{Sprint 3} & \thead{Sprint 4}
      & \thead{Tổng} \\
    \midrule
    \canDien & PO / BA & & & & & \\
    \canDien & SM / PM & & & & & \\
    \canDien & Dev Backend / STT & & & & & \\
    \canDien & Dev Frontend / Desktop & & & & & \\
    \canDien & Dev AI & & & & & \\
    \canDien & QC & & & & & \\
    \midrule
    \textbf{Tổng} & & & & & & \\
    \bottomrule
  \end{tabularx}
\end{table}

Khi đã có số liệu, biểu đồ tải theo sprint (resource histogram) được vẽ từ bảng này: cột là tổng giờ thực tế, chia
màu theo thành viên (biểu đồ cột chồng); đường ngang là năng lực của nhóm theo~\eqref{eq:capacity}. Sprint nào có
cột vượt đường năng lực là quá tải và phải san tải theo Mục~\ref{sec:leveling}.

\section{Phân công công việc theo sprint}

Bảng~\ref{tab:assignment} giao các công việc của lịch trình cho từng sprint. Mã công việc A--W như ở
Bảng~\ref{tab:cpm}.

\begin{table}[htbp]
  \centering
  \caption{Phân công công việc theo sprint}
  \label{tab:assignment}
  \small
  \begin{tabularx}{\linewidth}{@{}P{1.9cm}LP{3.4cm}@{}}
    \toprule
    \thead{Sprint} & \thead{Công việc} & \thead{Người phụ trách chính} \\
    \midrule
    1 & A Khởi động, B Kiến trúc, C Ghi âm, D Offline, E Upload, F STT & \canDien \\
    2 & G Live, H Benchmark LLM & \canDien \\
    3 & I--L Tóm tắt (\#68--\#71), M--Q Dịch (\#61--\#65), R Đánh giá (\#66, \#72), S Test E2E (\#99, \#100)
      & \canDien \\
    4 & T Desktop, U Streaming và STT remote, V Deploy AWS & \canDien \\
    Kết thúc & W UAT, phát hành, tài liệu & Cả nhóm \\
    \bottomrule
  \end{tabularx}
\end{table}

\section{San tải}
\label{sec:leveling}

Bốn tình huống quá tải đã được dự kiến, mỗi tình huống có cách xử lý dựa trên dự trữ của lịch trình hoặc biện pháp
cho rủi ro R14 (Bảng~\ref{tab:leveling}).

\begin{table}[htbp]
  \centering
  \caption{Các tình huống san tải}
  \label{tab:leveling}
  \small
  \begin{tabularx}{\linewidth}{@{}P{4.2cm}LP{3cm}@{}}
    \toprule
    \thead{Tình huống} & \thead{Cách xử lý} & \thead{Dựa trên} \\
    \midrule
    Sprint 3 dồn nhiều việc AI (tóm tắt và dịch song song) & Dev frontend làm giao diện \#71/\#64 sau khi API
      xong; QC viết test case song song & Dự trữ 3--5 ngày của K, L, O, P \\
    Dev STT quá tải ở Sprint 4 (G $\rightarrow$ U $\rightarrow$ V đều găng) & Giao T (desktop, dự trữ 4 ngày) cho
      Dev frontend; phần tài liệu của deploy (V) chia cho PM & Dự trữ của T \\
    Thành viên bận việc khác trong một giai đoạn & Dời việc không găng (E: 16 ngày dự trữ, J: 12 ngày) sang sau giai
      đoạn đó & CPM \\
    Một người nghỉ đột xuất & Code review chéo; README, ARCHITECTURE và HANDOFF để người khác tiếp quản & R14 \\
    \bottomrule
  \end{tabularx}
\end{table}

\section{Nguồn lực máy móc và dịch vụ}

Ngoài con người, dự án cần bảy loại nguồn lực máy móc và dịch vụ (Bảng~\ref{tab:machines}). Phần lớn đã có sẵn
hoặc miễn phí; các khoản phải trả tiền được tính ở Chương~\ref{chap:cost}.

\begin{table}[htbp]
  \centering
  \caption{Nguồn lực máy móc và dịch vụ}
  \label{tab:machines}
  \small
  \begin{tabularx}{\linewidth}{@{}P{3.4cm}LP{2.6cm}P{2.8cm}@{}}
    \toprule
    \thead{Nguồn lực} & \thead{Dùng cho} & \thead{Số lượng / cấu hình} & \thead{Chi phí} \\
    \midrule
    Laptop dev (i7-13620H, RTX 4050 6\,GB) & Phát triển, Whisper GPU local, build bộ cài & 1+ & Có sẵn \\
    Laptop thành viên khác (CPU) & Frontend, test, dùng engine remote & \canDien & Có sẵn \\
    AWS EC2 & Bản web test: \texttt{t3.large} (CPU) hoặc \texttt{g4dn.xlarge} (GPU, bật theo giờ) & 1
      & Mục~\ref{sec:nonlabour} \\
    AWS S3 & Audio, sao lưu database & 1 bucket & Vài nghìn VND/GB/tháng \\
    Whisper API (Groq/OpenAI) & STT cho bản web CPU & Theo giờ audio & Mục~\ref{sec:nonlabour} \\
    LLM API hoặc Ollama local & Tóm tắt, dịch & Theo token & Mục~\ref{sec:nonlabour} \\
    Azure Boards, GitHub & Quản lý công việc, mã nguồn & Gói miễn phí & 0 \\
    \bottomrule
  \end{tabularx}
\end{table}

% ============================================================ Chapter 7
\chapter{Chi phí và ngân sách}
\label{chap:cost}

\textbf{Chi phí chính là nhân công, tính theo đơn giá 100.000~VND/giờ trên dashboard chi phí của nhóm.} Chi phí
ngoài nhân công nhỏ: bản test trên AWS khoảng 2,5--3 triệu VND/tháng nếu chạy 24/7 với cấu hình CPU, và gần như
bằng 0 nếu chỉ bật khi demo. Ngân sách giữ thêm 10\% dự phòng.

\section{Đơn giá}

Hai đơn giá được dùng cho mọi phép tính chi phí trong tài liệu (Bảng~\ref{tab:rates}).

\begin{table}[htbp]
  \centering
  \caption{Đơn giá dùng trong kế hoạch chi phí}
  \label{tab:rates}
  \begin{tabularx}{\linewidth}{@{}LP{4cm}P{5cm}@{}}
    \toprule
    \thead{Hạng mục} & \thead{Đơn giá} & \thead{Nguồn} \\
    \midrule
    Nhân công (mọi vai trò) & 100.000 VND/giờ & Dashboard chi phí của nhóm \\
    Tỷ giá quy đổi & 26.000 VND/USD & Giả định \\
    \bottomrule
  \end{tabularx}
\end{table}

\section{Chi phí nhân công theo sprint}

Chi phí kế hoạch (PV) và chi phí thực tế (AC) của mỗi sprint bằng số giờ tương ứng nhân với đơn giá nhân công:
\begin{equation}
  \label{eq:labour-cost}
  \mathrm{PV}_{\text{sprint}} = \text{giờ kế hoạch} \times 100.000~\text{VND}, \qquad
  \mathrm{AC}_{\text{sprint}} = \text{giờ thực tế} \times 100.000~\text{VND}
\end{equation}
Tổng PV của các sprint là ngân sách nhân công khi hoàn thành (BAC). Các giá trị theo sprint
(Bảng~\ref{tab:labour}) là đầu vào cho phân tích EVA trong~\cite{mcr}.

\begin{table}[htbp]
  \centering
  \caption{Chi phí nhân công theo sprint}
  \label{tab:labour}
  \begin{tabularx}{\linewidth}{@{}LQ{2.6cm}Q{2.8cm}Q{2.6cm}Q{2.8cm}@{}}
    \toprule
    \thead{Sprint} & \thead{Giờ kế hoạch} & \thead{PV (VND)} & \thead{Giờ thực tế} & \thead{AC (VND)} \\
    \midrule
    Sprint 1 & \canDien & & \canDien & \\
    Sprint 2 & \canDien & & \canDien & \\
    Sprint 3 & \canDien & & \canDien & \\
    Sprint 4 & \canDien & & \canDien & \\
    \midrule
    \textbf{Tổng (BAC)} & & & & \\
    \bottomrule
  \end{tabularx}
\end{table}

\section{Chi phí ngoài nhân công}
\label{sec:nonlabour}

Giá trong Bảng~\ref{tab:nonlabour} là ước tính từ bảng giá công khai, cần kiểm tra lại trên AWS Pricing Calculator
và trang giá của nhà cung cấp trước khi dùng. Cấu hình A là bản web chạy 24/7 trên EC2 CPU, dùng Whisper API của
Groq và LLM cloud model nhỏ.

\begingroup\small
\setlength{\tabcolsep}{4pt}
\begin{longtable}{@{}P{3.6cm}P{3.6cm}P{3.2cm}P{4.4cm}@{}}
  \caption{Chi phí ngoài nhân công (ước tính mỗi tháng)}\label{tab:nonlabour}\\
  \toprule
  \thead{Khoản} & \thead{Cách tính} & \thead{Ước tính/tháng} & \thead{Ghi chú} \\
  \midrule
  \endfirsthead
  \multicolumn{4}{@{}l}{\footnotesize\itshape Bảng~\thetable{} (tiếp theo)}\\
  \toprule
  \thead{Khoản} & \thead{Cách tính} & \thead{Ước tính/tháng} & \thead{Ghi chú} \\
  \midrule
  \endhead
  \bottomrule
  \endlastfoot
  EC2 \texttt{t3.large} (Singapore), 24/7 & Khoảng 0,1 USD/giờ $\times$ 730 giờ & Khoảng 2,0 triệu VND
    & Bật/tắt theo giờ test thì rẻ hơn tương ứng \\
  EBS gp3 30\,GB & Khoảng 0,1 USD/GB & Khoảng 80.000 VND & -- \\
  Public IPv4 / Elastic IP & Khoảng 0,005 USD/giờ & Khoảng 95.000 VND & Nhớ release khi gỡ \\
  S3 Standard & Khoảng 0,025 USD/GB; 1 giờ audio khoảng 115\,MB & $<$ 50.000 VND cho 100 giờ audio & -- \\
  Whisper API -- Groq (\texttt{whisper-large-}\allowbreak\texttt{v3-turbo}) & Khoảng 0,04 USD/giờ audio
    & 100 giờ audio khoảng 104.000 VND & Live gửi lặp lại cửa sổ ngắn, nên chi phí thực tế cao hơn số giờ audio \\
  Whisper API -- OpenAI \texttt{whisper-1} & Khoảng 0,006 USD/phút & 100 giờ audio khoảng 940.000 VND
    & Phương án dự phòng \\
  LLM tóm tắt, dịch (cloud, model nhỏ) & Theo token; app log chi phí từng lần chạy & Khoảng 130.000 VND
    & Cấu hình \texttt{LLM\_PRICE\_*\_PER\_1M} để log chi phí \\
  EC2 GPU \texttt{g4dn.xlarge} (nếu cần) & Khoảng 0,7 USD/giờ & 20 giờ demo khoảng 370.000 VND
    & Không chạy 24/7 \\
  Chứng chỉ HTTPS, tên miền & Let's Encrypt và sslip.io miễn phí & 0 & Tên miền riêng tuỳ chọn \\
  Ký số bộ cài (code signing) & -- & 0 & Ngoài phạm vi (R12) \\
  \midrule
  \textbf{Tổng cấu hình A, 24/7} & & \textbf{Khoảng 2,5--3 triệu VND} & \\
\end{longtable}
\endgroup

Hình~\ref{fig:nonlabour} cho thấy cơ cấu của cấu hình A: máy chủ EC2 chiếm khoảng bốn phần năm chi phí, còn các API
STT và LLM cộng lại chưa tới 10\%. Vì vậy, cách tiết kiệm hiệu quả nhất là tắt máy khi không test, không phải cắt
giảm tính năng AI.

\begin{figure}[htbp]
  \centering
  \begin{tikzpicture}
    \begin{axis}[
        width=0.8\linewidth, height=6cm,
        xbar, bar width=0.42cm,
        xmin=0, xmax=2450, xtick distance=500,
        symbolic y coords={EC2,EBS,IP,S3,STT,LLM}, ytick=data, y dir=reverse,
        yticklabels={{EC2 t3.large},{EBS gp3 30\,GB},{Public IPv4},{S3 (100 giờ audio)},
                     {Whisper API (Groq)},{LLM API}},
        xlabel={Nghìn VND/tháng (ước tính)},
        nodes near coords, every node near coord/.append style={font=\footnotesize},
        grid=major, ymajorgrids=false, grid style={gray!25},
        tick label style={font=\small}, enlarge y limits=0.12,
        /pgf/number format/.cd, use comma, 1000 sep={.}]
      \addplot[fill=primaryblue, draw=primaryblue]
        coordinates {(2000,EC2) (80,EBS) (95,IP) (50,S3) (104,STT) (130,LLM)};
    \end{axis}
  \end{tikzpicture}
  \caption{Cơ cấu chi phí ngoài nhân công mỗi tháng của cấu hình A (S3 lấy mức tối đa 50 nghìn VND)}
  \label{fig:nonlabour}
\end{figure}

\section{Tổng ngân sách}

Tổng ngân sách gồm nhân công, hạ tầng và API trong thời gian test, cộng 10\% dự phòng
(Bảng~\ref{tab:total-budget}). Phần hạ tầng giả định thời gian test là 2 tháng với cấu hình A.

\begin{table}[htbp]
  \centering
  \caption{Tổng ngân sách}
  \label{tab:total-budget}
  \begin{tabularx}{\linewidth}{@{}LP{7cm}@{}}
    \toprule
    \thead{Khoản} & \thead{Giá trị} \\
    \midrule
    Nhân công (BAC nhân công) & \canDien[tổng giờ $\times$ 100.000] \\
    Hạ tầng và API trong thời gian test (giả định 2 tháng) & Khoảng 6 triệu VND \\
    Dự phòng & 10\% $\times$ (nhân công + hạ tầng) \\
    \midrule
    \textbf{Tổng ngân sách được duyệt} & \canDien \\
    \bottomrule
  \end{tabularx}
\end{table}

\section{Kiểm soát chi phí}

\begin{itemize}
  \item \textbf{Theo dõi:} dashboard chi phí của nhóm cập nhật giờ mỗi sprint; CPI và SPI được tính trong báo cáo
        trạng thái và phân tích EVA~\cite{mcr}.
  \item \textbf{Cảnh báo:} AWS Budget gửi email ở mức 80\%; log chi phí token của từng lần tóm tắt và dịch.
  \item \textbf{Ngưỡng hành động:} khi CPI $<$ 0,9 hoặc chi phí ngoài nhân công vượt 25\% ngân sách tháng, PM báo PO
        và Sponsor, cân nhắc cắt phạm vi (story Should, Could) hoặc chuyển sang cấu hình rẻ hơn (CPU kết hợp API,
        tắt máy khi không dùng).
  \item \textbf{Biện pháp kỹ thuật đã có:} \texttt{AI\_MOCK=true} khi demo; VAD không gửi đoạn im lặng;
        \texttt{LIVE\_INTERVAL\_SEC} giảm số request khi ghi live; engine remote dùng chung một server GPU cho nhiều
        máy khách.
\end{itemize}

% ============================================================ Chapter 8
\chapter{Tổ chức nhóm và truyền thông}
\label{chap:team}

\textbf{Nhóm làm Scrum với một Product Owner, một Scrum Master kiêm Project Manager, BA, QC và các Developer; mỗi
sản phẩm bàn giao có đúng một người chịu trách nhiệm cuối (A).} Kênh trao đổi chính là Azure Boards (công việc),
GitHub (mã nguồn) và nhóm chat (trao đổi nhanh).

\section{Thành viên và vai trò}

Bảng~\ref{tab:roles} gán vai trò Scrum và vai trò chuyên môn cho từng thành viên. Một người có thể giữ nhiều vai;
\canDien[chỉnh theo phân vai thật của nhóm].

\begin{table}[htbp]
  \centering
  \caption{Thành viên và vai trò}
  \label{tab:roles}
  \small
  \begin{tabularx}{\linewidth}{@{}P{2.6cm}P{3.2cm}P{3cm}L@{}}
    \toprule
    \thead{Thành viên} & \thead{Vai trò Scrum} & \thead{Chuyên môn} & \thead{Mảng phụ trách chính} \\
    \midrule
    \canDien & Product Owner & BA & Backlog, acceptance criteria, mẫu tham chiếu để đánh giá \\
    \canDien & Scrum Master / PM & Dev & Kế hoạch, theo dõi chi phí (CPI/SPI), rủi ro, báo cáo \\
    \canDien & Developer & Backend / STT & FastAPI, Whisper, pipeline, deploy \\
    \canDien & Developer & Frontend / Desktop & React, recorder, offline, Electron \\
    \canDien & Developer & AI (LLM) & Tóm tắt, dịch, đánh giá chất lượng \\
    \canDien & Developer & QC & Test case, Playwright, báo lỗi \\
    \bottomrule
  \end{tabularx}
\end{table}

\section{Ma trận RACI}
\label{sec:raci}

Ma trận RACI (Bảng~\ref{tab:raci}) dùng bốn ký hiệu: \textbf{R} -- thực hiện, \textbf{A} -- chịu trách nhiệm cuối
(mỗi dòng đúng một A), \textbf{C} -- được hỏi ý kiến, \textbf{I} -- được thông báo. Trong bảng, BE, FE là viết tắt
của backend và frontend.

\begin{table}[htbp]
  \centering
  \caption{Ma trận RACI}
  \label{tab:raci}
  \small
  \setlength{\tabcolsep}{3pt}
  \begin{tabularx}{\linewidth}{@{}LQ{0.95cm}Q{1.35cm}Q{1.3cm}Q{1.3cm}Q{1.3cm}Q{0.95cm}Q{1.55cm}@{}}
    \toprule
    \thead{Công việc / sản phẩm} & \thead{PO} & \thead{PM/\allowbreak SM} & \thead{Dev BE} & \thead{Dev FE} & \thead{Dev AI}
      & \thead{QC} & \thead{Sponsor} \\
    \midrule
    Project Charter, Business Case & C & \textbf{A/R} & I & I & I & I & C (duyệt) \\
    Product Backlog, user story, acceptance criteria & \textbf{A/R} & C & C & C & C & C & I \\
    Ưu tiên backlog, chấp nhận story & \textbf{A} & C & I & I & I & C & I \\
    Kế hoạch sprint, lịch, ước lượng & C & \textbf{A/R} & R & R & R & R & I \\
    Kiến trúc, lựa chọn công nghệ & C & C & \textbf{A/R} & R & R & I & I \\
    Ghi âm, offline, upload (E1) & C & I & R & \textbf{A/R} & I & C & I \\
    STT, live streaming, engine remote & C & I & \textbf{A/R} & R & I & C & I \\
    Tóm tắt, dịch (\#67, \#60) & C & I & R & R & \textbf{A/R} & C & I \\
    App desktop & C & I & R & \textbf{A/R} & I & C & I \\
    Deploy AWS, vận hành bản test & I & C & \textbf{A/R} & I & I & C & I \\
    Test case, test E2E, báo cáo test & C & I & C & C & C & \textbf{A/R} & I \\
    Definition of Done, nghiệm thu story & \textbf{A} & C & R & R & R & R & I \\
    Quản lý rủi ro, Change Log & C & \textbf{A/R} & C & C & C & C & I \\
    Theo dõi chi phí (EVA, CPI/SPI) & I & \textbf{A/R} & I & I & I & I & I \\
    Báo cáo trạng thái, sprint report & C & \textbf{A/R} & C & C & C & C & I \\
    Nghiệm thu cuối, SOW & R & R & I & I & I & C & \textbf{A} \\
    \bottomrule
  \end{tabularx}
\end{table}

PM/SM giữ vai A ở năm trên mười sáu dòng, đều là các sản phẩm quản lý; PO giữ vai A ở các quyết định về sản phẩm
(backlog, ưu tiên, Definition of Done). Sponsor chỉ giữ vai A ở bước nghiệm thu cuối, đúng với phân quyền trong
điều lệ dự án~\cite{pid}.

\section{Lịch họp Scrum}

Năm sự kiện Scrum được tổ chức theo Bảng~\ref{tab:meetings}; mỗi sự kiện có đầu ra cụ thể để kiểm tra.

\begin{table}[htbp]
  \centering
  \caption{Lịch họp Scrum}
  \label{tab:meetings}
  \small
  \begin{tabularx}{\linewidth}{@{}P{2.6cm}P{2.4cm}Q{1.6cm}P{3.2cm}L@{}}
    \toprule
    \thead{Sự kiện} & \thead{Khi nào} & \thead{Thời lượng} & \thead{Thành phần} & \thead{Đầu ra} \\
    \midrule
    Sprint Planning & Đầu sprint & 1--2 giờ & Cả nhóm & Sprint backlog, mục tiêu sprint, ước lượng (Planning
      Poker) \\
    Daily Scrum & \canDien[hằng ngày hoặc 3 buổi/tuần] & 15 phút & Dev, SM (PO tuỳ chọn) & Đã làm, sẽ làm,
      vướng mắc \\
    Backlog Refinement & Giữa sprint & 1 giờ & PO, BA, đại diện Dev/QC & Story sẵn sàng cho sprint sau (có
      acceptance criteria, đã ước lượng) \\
    Sprint Review & Cuối sprint & 1 giờ & Cả nhóm, Sponsor, người dùng thử (khi cần) & Demo, phản hồi, cập nhật
      backlog \\
    Retrospective & Sau Sprint Review & 45 phút & Cả nhóm & 1--3 hành động cải tiến~\cite{mcr} \\
    \bottomrule
  \end{tabularx}
\end{table}

\section{Kế hoạch truyền thông}

Kế hoạch truyền thông (Bảng~\ref{tab:comms}) xác định ai nhận thông tin gì, qua kênh nào và với tần suất nào.

\begin{table}[htbp]
  \centering
  \caption{Kế hoạch truyền thông}
  \label{tab:comms}
  \small
  \setlength{\tabcolsep}{4pt}
  \begin{tabularx}{\linewidth}{@{}P{3.3cm}P{2.6cm}LP{2.1cm}P{2.1cm}@{}}
    \toprule
    \thead{Thông tin} & \thead{Gửi cho} & \thead{Hình thức} & \thead{Tần suất} & \thead{Người gửi} \\
    \midrule
    Tiến độ task, lỗi & Cả nhóm & Azure Boards (board, Bug) & Liên tục & Người làm task \\
    Vướng mắc cần gỡ & SM, người liên quan & Nhóm chat, nêu trong Daily & Khi phát sinh & Bất kỳ \\
    Code, review & Dev & GitHub Pull Request trên nhánh feature & Mỗi task & Dev \\
    Báo cáo trạng thái (burndown, CPI/SPI, rủi ro) & Sponsor, cả nhóm & Sprint report~\cite{mcr}, dashboard chi phí
      & Mỗi sprint & PM \\
    Demo tính năng & Sponsor, người dùng thử & Sprint Review, bản web trên AWS & Mỗi sprint & PO \\
    Thay đổi phạm vi & Sponsor, cả nhóm & Change Log~\cite{mcr}, email hoặc chat & Khi có thay đổi & PO/PM \\
    Sự cố bản deploy & Dev ops, PM & Nhóm chat & Ngay & Người phát hiện \\
    \bottomrule
  \end{tabularx}
\end{table}

Ba quy tắc trao đổi đi kèm: câu hỏi kỹ thuật được trả lời trong ngày làm việc; quyết định về phạm vi phải ghi lại
trên board hoặc Change Log, không để lại chỉ trong chat; tài liệu dùng chung nằm trong repo (thư mục
\path{docs/}).

\section{Leo thang vấn đề}

Một vấn đề được xử lý ở cấp thấp nhất có thể, và chỉ chuyển lên cấp trên khi vượt quyền hoặc vượt thời gian của
cấp dưới:
\begin{enumerate}
  \item Người làm task tự xử lý (trong tối đa 1 ngày).
  \item Báo Scrum Master trong Daily; SM điều phối người hỗ trợ.
  \item Nếu ảnh hưởng mục tiêu sprint: PO quyết định cắt hoặc dời story.
  \item Nếu ảnh hưởng mốc, ngân sách vượt 10\% hoặc phạm vi dự án: PM báo Sponsor để xin phê duyệt.
\end{enumerate}

% ============================================================ Chapter 9
\chapter{Kết luận và phê duyệt}
\label{chap:conclusion}

\section{Tóm tắt}

Kế hoạch đã xác định cách thực hiện và kiểm soát dự án SumTalk:
\begin{itemize}
  \item \textbf{Phạm vi:} WBS gồm 9 nhánh cấp 1 và 33 gói công việc, mỗi gói truy vết được tới task trên Azure
        Boards hoặc file trong repo.
  \item \textbf{Effort:} 118 SP đề xuất cho 18 hạng mục; Function Point (khoảng 1.560 giờ) và COCOMO~81
        (3.500--5.100 giờ) xác nhận quy mô sản phẩm.
  \item \textbf{Lịch trình:} 23 công việc, đường găng 33 ngày làm việc qua ghi âm, STT, live, streaming và deploy;
        nhánh tóm tắt và dịch có 3--12 ngày dự trữ.
  \item \textbf{Rủi ro:} 14 rủi ro, 5 ở vùng Cao; R2 đã giảm xuống mức Thấp nhờ engine remote, R11 đang xảy ra.
  \item \textbf{Nguồn lực và chi phí:} ràng buộc là giờ rảnh của từng người; đơn giá nhân công là 100.000 VND mỗi
        giờ, hạ tầng và API khoảng 2,5--3 triệu VND mỗi tháng, dự phòng 10\%.
  \item \textbf{Tổ chức:} ma trận RACI với một người A cho mỗi sản phẩm, năm sự kiện Scrum, kế hoạch truyền thông và
        bốn cấp leo thang.
\end{itemize}

\section{Việc cần làm trước khi trình duyệt}

\begin{enumerate}
  \item Điền các ô \canDien: người lập, tên thành viên theo vai trò, người phụ trách theo sprint, giờ kế hoạch và
        giờ thực tế theo người và theo sprint, velocity của Sprint 1--3, lịch Daily và tổng ngân sách được duyệt.
  \item Thay ngày bắt đầu giả định 07/09/2026 bằng ngày bắt đầu Sprint 1 thật, rồi tính lại lịch bằng
        \texttt{pm\_calc.py}~\cite{pmcalc}.
  \item Xác nhận tổng story point của Bảng~\ref{tab:estimate} (118 SP, phiên bản 2.0 ghi 125 SP) với điểm đã vote
        trên Azure Boards.
  \item Đối chiếu R1--R5 với biên bản họp, và đối chiếu release plan với việc đã thực sự làm trong từng sprint.
  \item Thay vai trò in nghiêng ở cột phụ trách của Risk Register bằng tên thành viên.
\end{enumerate}

\needspace{14\baselineskip}
\section{Phê duyệt}

Nhóm dự án đề nghị Sponsor, Product Owner và Project Manager phê duyệt kế hoạch này làm đường cơ sở về phạm vi,
lịch trình và chi phí để theo dõi và kiểm soát dự án SumTalk.

\begin{table}[H]
  \centering
  \renewcommand{\arraystretch}{2.4}
  \begin{tabularx}{\linewidth}{@{}P{3.6cm}LP{3.2cm}P{2.6cm}@{}}
    \toprule
    \thead{Vai trò} & \thead{Họ tên} & \thead{Chữ ký} & \thead{Ngày} \\
    \midrule
    Sponsor & & & \\
    Product Owner & & & \\
    Project Manager & & & \\
    \bottomrule
  \end{tabularx}
\end{table}

% ============================================================ References
\begin{thebibliography}{9}
\addcontentsline{toc}{chapter}{\bibname}
  \bibitem{pid} Nhóm SumTalk, \emph{Hồ sơ khởi động dự án (Project Initiation Document)}, Tài liệu 01,
    mã SUMTALK-PID, phiên bản 2.1, 24/11/2026.
  \bibitem{srs} Nhóm SumTalk, \emph{Đặc tả yêu cầu phần mềm (Software Requirements Specification)},
    Tài liệu 02, mã SUMTALK-SRS, phiên bản 2.1, 24/11/2026.
  \bibitem{qtd} Nhóm SumTalk, \emph{Chất lượng và kiểm thử (Quality \& Test Document)}, Tài liệu 04,
    mã SUMTALK-QTD, phiên bản 2.0, 13/11/2026.
  \bibitem{mcr} Nhóm SumTalk, \emph{Báo cáo theo dõi và kiểm soát (Monitoring \& Control Report)}, Tài liệu 05,
    mã SUMTALK-MCR, phiên bản 2.0, 13/11/2026.
  \bibitem{pcr} Nhóm SumTalk, \emph{Báo cáo kết thúc dự án (Project Closure Report)}, Tài liệu 06,
    mã SUMTALK-PCR, phiên bản 2.0, 13/11/2026.
  \bibitem{pmcalc} Nhóm SumTalk, \emph{Công cụ tính COCOMO, NPV, ROI, thời gian hoàn vốn và CPM},
    kho mã nguồn \texttt{DatTapTanhCode/SumTalk}, tệp \path{docs/project/tools/pm_calc.py}.
  \bibitem{pmbok} Project Management Institute, \emph{A Guide to the Project Management Body of Knowledge
    (PMBOK Guide)}, 6th ed., Newtown Square, PA: PMI, 2017.
  \bibitem{boehm} B. W. Boehm, \emph{Software Engineering Economics}, Englewood Cliffs, NJ: Prentice-Hall, 1981.
  \bibitem{ifpug} International Function Point Users Group, \emph{Function Point Counting Practices Manual},
    Release 4.3.1, 2010.
\end{thebibliography}

\end{document}
